\chapter{System Analysis and Design}
\label{ch:analysis-design}

\section{Introduction}

This chapter presents the system analysis and design of the EcoTrace: Intelligent
Electronic Waste Recovery and Circular Economy Management System. It
explains the methodology adopted during the development process, examines the
existing electronic waste management practices, identifies their major
weaknesses, and describes the proposed system developed to address those
weaknesses.

The chapter also presents the functional and non-functional requirements of
EcoTrace, the high-level system architecture, use case model, data flow and
sequence models, database design, and user interface design. Furthermore, the
programming languages, frameworks, cloud services, development tools, and
platform requirements used in the development of the system are discussed.

The analysis and design presented in this chapter provide the technical
foundation for the implementation and testing activities described in
Chapter~\ref{ch:implementation}.

\section{Methodology}

\subsection{Adopted Software Development Methodology}

The development of EcoTrace followed an adapted Agile Software Development
Methodology. Agile development is an iterative and incremental approach in
which a software system is developed through a series of manageable cycles
rather than being implemented as a single large process. It emphasizes continuous
improvement, frequent testing, responsiveness to changing requirements, and the
delivery of usable software components throughout the development lifecycle
\citep{agile2001,sommerville2016}.

The Agile methodology was selected because EcoTrace consists of several
interdependent modules, including authentication, electronic waste registration,
artificial intelligence classification, pickup scheduling, reverse logistics, repair,
refurbishment, recycling, notifications, mapping, marketplace management,
reporting, and administration. Developing these modules incrementally made it
possible to implement, test, review, and improve each component before
integrating it into the complete system.

The methodology was adapted to suit an individual final-year software
development project. Therefore, the project did not implement every formal
Scrum role or ceremony. Instead, Agile principles were applied through iterative
planning, incremental implementation, continuous testing, regular review, and
progressive integration of the system modules.

\subsection{Justification for the Agile Methodology}

The Agile methodology was considered suitable for EcoTrace for several reasons.

\begin{enumerate}[leftmargin=1.4cm]
  \item \textbf{Support for changing requirements}: The requirements of EcoTrace
    evolved during development as additional services such as Render
    deployment, Cloudinary image hosting, Firebase Cloud Messaging, and route
    optimization were integrated. Agile enabled these changes to be incorporated
    without restarting the entire development process.
  \item \textbf{Incremental module development}: The system contains multiple
    modules that could be designed, implemented, and tested independently
    before being integrated into the complete platform.
  \item \textbf{Early identification of defects}: Continuous testing during each
    development iteration helped identify interface, authentication, API,
    database, and deployment problems before they affected the complete
    system.
  \item \textbf{Improved risk management}: Technically complex components,
    including cloud integration, image upload, location services, and artificial
    intelligence, could be evaluated progressively rather than being postponed
    until the final stage.
  \item \textbf{Cross-platform development support}: The Android, Web, and
    Windows applications could be developed and reviewed incrementally while
    sharing a common Flutter codebase.
  \item \textbf{Continuous improvement}: Feedback obtained during
    implementation and testing could be used to refine the user interfaces,
    workflows, database structure, and backend services.
  \item \textbf{Faster delivery of functional components}: Usable modules such
    as registration, login, waste submission, image upload, and pickup
    scheduling could be completed before every advanced feature was finalized.
\end{enumerate}

\subsection{Application of Agile to the EcoTrace Project}

The adapted Agile development process used for EcoTrace consisted of the
following iterative activities.

\textbf{Requirements Identification and Planning.} The first activity involved
identifying the electronic waste management problem, reviewing existing
systems, defining the project objectives, and determining the major system users
and functional modules. The identified user categories included household users,
business users, collectors, drivers, technicians, recyclers, environmental officers,
administrators, and super administrators. Their expected responsibilities were
analysed and translated into functional and non-functional requirements. The
requirements were prioritized according to their importance, technical
complexity, dependencies, and contribution to the main electronic waste
management workflow.

\textbf{System Analysis.} Existing electronic waste management practices were
examined to understand how electronic waste is generated, reported, collected,
transported, repaired, refurbished, recycled, and monitored. The analysis
identified weaknesses such as manual documentation, fragmented
communication, poor lifecycle traceability, inefficient collection planning, limited
environmental reporting, and the absence of an integrated cross-platform system.
The results of the analysis were used to define the responsibilities, workflows,
data requirements, and system boundaries of EcoTrace.

\textbf{System Design.} The design activity translated the identified
requirements into technical models and implementation specifications. The
system architecture, user roles, use cases, data flows, database collections,
interfaces, and external service integrations were designed before the
corresponding modules were implemented. EcoTrace was designed as a
cross-platform client-server system. Flutter provides the Android, Web, and
Windows interfaces, while FastAPI provides the central backend services. Cloud
Firestore stores structured system data, Cloudinary hosts uploaded images,
Firebase Authentication manages user identity, Firebase Cloud Messaging
provides notifications, Google Maps Platform provides location and routing
services, and PyTorch supports electronic waste image classification.

\textbf{Incremental Implementation.} Implementation was divided into
manageable iterations. Core services were implemented before dependent and
advanced modules. The initial iterations focused on project configuration,
authentication, user profiles, role management, and navigation. Subsequent
iterations implemented electronic waste registration, image upload, Cloud
Firestore integration, pickup requests, location services, and notification
functions. Additional iterations addressed reverse logistics, repair, refurbishment,
recycling, reporting, administrative dashboards, Render deployment, and
artificial intelligence integration. This incremental approach reduced development
complexity and enabled completed features to be tested before additional
modules were added.

\textbf{Testing and Validation.} Testing was performed continuously throughout
development rather than being postponed until the entire system was complete.
Individual functions and interface components were tested during unit testing.
Interactions between Flutter, FastAPI, Firebase, Cloudinary, Google Maps
Platform, Render, and PyTorch were examined during integration testing. Complete workflows were
subsequently evaluated through system testing. The results of each test cycle were
used to correct defects, improve validation rules, refine interfaces, and strengthen
system reliability. Detailed testing procedures and results are presented in
Chapter~\ref{ch:implementation}.

\textbf{Review and Refinement.} At the end of each iteration, the implemented
features were reviewed against the corresponding requirements. Incomplete,
inconsistent, or inefficient functions were revised before the next iteration began.
The review process also considered user-interface consistency, data integrity,
authentication security, API responses, navigation, cross-platform behaviour, and
the accuracy of information displayed to users.

\textbf{Deployment.} The backend was deployed on the Render Free Web Service
after its main endpoints and external integrations were configured. Sensitive
credentials were managed through environment variables rather than being
stored directly in the source code. Cloud Firestore, Firebase Authentication,
Firebase Cloud Messaging, Cloudinary, and Google Maps Platform remained
externally managed cloud services integrated with the deployed backend and
Flutter applications. The
free Render service was suitable for academic demonstration and testing, although
its cold-start behaviour and limited computing resources were recognized as
deployment constraints.

\textbf{Maintenance and Future Improvement.} The iterative methodology allows
EcoTrace to be maintained and expanded after the initial implementation. New
modules, user roles, AI models, reports, integrations, and environmental
indicators can be introduced without redesigning the complete system. Future
development iterations may improve offline functionality, large-scale deployment,
advanced route optimization, payment integration, IoT sensor support, and the
accuracy and efficiency of the electronic waste classification model.
Figure~\ref{fig:agile-methodology} illustrates the
iterative development process adopted for the design, implementation, testing,
and refinement of the EcoTrace system.

\begin{figure}[H]
  \centering
  \imageplaceholder{Adapted Agile Software Development Methodology used for EcoTrace}
  \caption{Adapted Agile Software Development Methodology used for EcoTrace. Source: Author's design based on Agile software development principles \citep{agile2001,sommerville2016}.}
  \label{fig:agile-methodology}
\end{figure}

\section{Analysis of the Existing System}

The existing electronic waste management process in Sierra Leone is largely
manual, fragmented, and informal. Electronic waste is generated by households,
businesses, educational institutions, healthcare facilities, government offices,
repair centres, and other organizations. However, there is no widely adopted
integrated digital platform for registering, collecting, transporting, repairing,
refurbishing, recycling, and monitoring electronic waste throughout its lifecycle.

In the existing process, individuals and organizations commonly store
obsolete electronic devices, dispose of them together with general household
waste, sell them to informal collectors, donate them, or deliver them to repair
technicians. Devices that cannot be repaired may be dismantled manually to
recover selected components or materials. Residual waste may then be discarded
through uncontrolled channels.

The major activities within the existing system are described in the following
subsections.

\subsection{Electronic Waste Generation}
Electronic waste is generated when electrical or electronic equipment becomes
obsolete, damaged, unwanted, or uneconomical to repair. Common examples
include mobile phones, laptops, desktop computers, televisions, printers,
batteries, household appliances, cables, circuit boards, and accessories.
Households and organizations often keep obsolete devices for long periods
because they lack reliable information about authorized collection centres,
recycling facilities, or environmentally safe disposal methods. Consequently,
significant quantities of electronic waste remain unreported and untracked.

\subsection{Waste Identification and Classification}
Electronic waste is generally identified manually by its owner, collector,
technician, or recycler. Classification depends on personal knowledge and visual
inspection rather than a standardized digital process. There is limited support for
automatically identifying electronic waste categories from images. Incorrect or
incomplete classification can affect collection planning, repair decisions,
recycling procedures, and the accuracy of environmental reports.

\subsection{Collection Requests}
The existing process does not provide a centralized digital mechanism through
which households and businesses can register electronic waste and request
collection. Collection arrangements may be made through telephone calls,
personal contacts, informal agreements, or direct delivery to repairers and
recyclers. These methods provide limited documentation and make it difficult to
confirm when a request was submitted, accepted, assigned, collected, or
completed.

\subsection{Pickup Scheduling}
Pickup schedules are commonly arranged manually between waste owners and
collectors. Information such as the waste category, quantity, pickup location,
preferred date, collector assignment, and current status may not be recorded in a
structured format. Manual scheduling increases the possibility of missed
appointments, delayed collections, duplicate assignments, and communication
failures.

\subsection{Collection and Transportation}
Collectors and transport personnel may receive collection information through
telephone calls, written notes, or verbal instructions. Routes are often selected
based on personal experience rather than systematic distance calculations or
route optimization. The absence of integrated location and routing services may
result in longer travel distances, inefficient fuel consumption, delayed pickups,
and limited visibility into transportation activities. There is also limited digital
evidence showing when an item was collected, who collected it, which vehicle
transported it, and the facility to which it was delivered.

\subsection{Repair and Refurbishment}
Repair technicians assess devices manually to determine whether they can be
restored to working condition. Repair records, replaced components, estimated
costs, diagnostic findings, and final outcomes may be kept in personal notebooks
or not documented at all. Refurbished devices may be returned to their owners,
sold informally, or stored for future use. The absence of centralized records limits
traceability and makes it difficult to determine how many devices were
successfully repaired, refurbished, reused, or rejected.

\subsection{Recycling and Material Recovery}
Devices that cannot be repaired may be dismantled to recover reusable
components and valuable materials. In informal environments, dismantling may
involve unsafe manual techniques, open burning, or uncontrolled disposal of
hazardous components. Recycling records are often incomplete or unavailable.
Therefore, it is difficult to determine the quantity of electronic waste received,
the materials recovered, the residual waste produced, or the final disposal
method used.

\subsection{Marketplace and Redistribution}
Reusable devices, repaired equipment, spare parts, and recovered materials may
be sold through informal markets, social networks, repair shops, or direct
personal contacts. The existing process does not provide a centralized
marketplace linked to verified electronic waste recovery records. As a result,
buyers may have limited information about the condition, origin, inspection
status, or refurbishment history of an item.

\subsection{Communication and Notifications}
Communication among waste owners, collectors, drivers, technicians, recyclers,
and administrators is mainly conducted through telephone calls, text messages,
or face-to-face interaction. These communication methods are not integrated with
operational records. Users may therefore fail to receive timely information about
pickup approval, collector assignment, schedule changes, transportation status,
repair progress, recycling completion, or system announcements.

\subsection{Record Keeping and Reporting}
Existing records are commonly stored in notebooks, spreadsheets, individual files,
or separate organizational systems. Such records may be incomplete, duplicated,
inaccessible, or difficult to verify. The lack of centralized data prevents
administrators and environmental officers from generating accurate reports on
the amount and categories of electronic waste generated, collection requests and
completed pickups, collection locations and operational coverage, repaired and
refurbished devices, recycled materials and recovery outcomes, processing
facilities and responsible personnel, environmental performance indicators, and
unresolved or delayed electronic waste cases. Without reliable data, planning,
monitoring, compliance assessment, and policy development become difficult.

\subsection{Environmental Monitoring}
Environmental officers require accurate and timely information to monitor
electronic waste collection, transportation, processing, and disposal. Under the
existing process, monitoring often depends on physical inspections, manually
submitted records, and reports from individual organizations. This fragmented
approach limits real-time visibility and makes it difficult to identify collection
gaps, unsafe processing activities, waste hotspots, delayed cases, and trends in
electronic waste generation.

\subsection{Existing System Workflow}
The general workflow of the existing electronic waste management process can
be summarized as follows:

\begin{enumerate}[leftmargin=1.4cm]
  \item A household or organization identifies an obsolete or damaged electronic
    device.
  \item The owner stores the device, disposes of it with general waste, contacts an
    informal collector, or takes it to a technician.
  \item The collector or technician visually examines the device.
  \item Collection and transportation arrangements are made manually.
  \item Repairable devices are repaired or resold.
  \item Non-repairable devices are dismantled for components or materials.
  \item Remaining waste is recycled, stored, burned, dumped, or otherwise
    discarded.
  \item Records, where available, are maintained separately by individual
    participants.
\end{enumerate}

This workflow lacks a centralized mechanism for registration, classification,
assignment, tracking, reporting, and environmental monitoring.
Figure~\ref{fig:existing-workflow} illustrates the fragmented workflow associated
with the existing electronic waste management process.

\begin{figure}[H]
  \centering
  \imageplaceholder{Existing Electronic Waste Management Workflow}
  \caption{Existing Electronic Waste Management Workflow. Source: Author's analysis of the existing process.}
  \label{fig:existing-workflow}
\end{figure}

\section{Problems of the Existing System}

The analysis of the existing electronic waste management process revealed
several operational, technical, environmental, and administrative weaknesses.
These problems reduce the efficiency, traceability, safety, and sustainability of
electronic waste management activities. The major problems identified are
discussed in the following subsections.

\subsection{Lack of a Centralized Information System}
The existing process does not provide a centralized platform for registering,
storing, updating, and retrieving electronic waste information. Records may be
maintained separately by households, businesses, collectors, technicians,
recyclers, and public institutions. This fragmentation makes it difficult to obtain a
complete and consistent view of electronic waste activities. The absence of
centralized data also increases the likelihood of duplicated records, missing
information, inconsistent waste status updates, inaccurate reporting, and poor
coordination among stakeholders.

\subsection{Manual and Unstructured Record Keeping}
Electronic waste records are frequently maintained using notebooks, paper forms,
spreadsheets, telephone messages, or informal verbal communication. Manual
record keeping is vulnerable to human error, loss, damage, unauthorized
alteration, and incomplete documentation. It also makes searching, updating,
validating, and analysing records difficult. As the volume of electronic waste
increases, manual documentation becomes increasingly inefficient and unsuitable
for large-scale management.

\subsection{Poor Electronic Waste Traceability}
The existing system provides limited ability to track an electronic waste item
from the point of generation to its final outcome. It may be difficult to determine
who registered or owned the item, when and where it was collected, which
collector or driver handled it, the facility to which it was transported, whether it
was repaired, refurbished, recycled, or disposed of, which materials were
recovered, and whether the item was processed safely. This lack of traceability
reduces transparency and accountability across the electronic waste lifecycle.

\subsection{Absence of Standardized Waste Classification}
Electronic waste is commonly classified through manual visual inspection. The
accuracy of this process depends on the knowledge and experience of the
individual performing the assessment. Different users may describe similar
devices using different names or categories. Such inconsistencies affect collection
planning, storage, repair decisions, recycling procedures, data analysis, and
reporting. The absence of an AI-assisted classification tool also places the full
burden of identification on users and waste-management personnel.

\subsection{Inefficient Pickup Scheduling}
Pickup requests and schedules are commonly arranged through telephone calls,
personal contacts, or informal agreements. This approach does not provide a
reliable mechanism for recording the date and time of the request, the preferred
collection date, the waste category and quantity, the pickup location, the assigned
collector or driver, the current request status, and reasons for cancellation or
delay. As a result, pickups may be delayed, forgotten, duplicated, or assigned
inefficiently.

\subsection{Limited Route Planning and Optimization}
Collectors and drivers often select collection routes using personal judgement
rather than calculated distance, travel time, traffic conditions, and the location of
multiple pickup requests. The absence of route optimization may result in longer
travel distances, higher fuel consumption, increased transportation costs, delayed
pickups, inefficient vehicle utilization, and increased environmental emissions.
There is also limited real-time visibility into collector and driver activities.

\subsection{Weak Coordination Among Stakeholders}
Households, businesses, collectors, drivers, technicians, recyclers, environmental
officers, and administrators often operate independently. Information is
exchanged through telephone calls, messages, and verbal communication rather
than through a shared operational system. This weak coordination may cause
delays, misunderstandings, incomplete handover records, duplicated activities,
and uncertainty regarding responsibility for specific electronic waste items.

\subsection{Inadequate Communication and Notification Mechanisms}
The existing process does not provide automated notifications for important
events. Users may not receive timely information when a pickup request is
submitted, a request is approved or rejected, a collector is assigned, a pickup date
changes, an item is collected, repair or refurbishment is completed, recycling is
completed, or an environmental announcement is issued. This communication
gap reduces user confidence and operational responsiveness.

\subsection{Poor Repair and Refurbishment Documentation}
Repair and refurbishment activities are not consistently recorded in a centralized
system. Important information may be missing, including diagnostic findings,
device condition, replaced components, repair cost, assigned technician, repair
status, quality-assurance results, and final refurbishment outcome. The lack of
documentation limits accountability and makes it difficult to verify the condition
and history of refurbished devices.

\subsection{Limited Recycling and Material-Recovery Records}
Recycling activities may not be documented in sufficient detail. The existing
process provides limited information about the quantity of electronic waste
received, categories of devices processed, recovered metals, plastics, glass, and
components, hazardous materials identified, residual waste generated,
responsible recycling personnel, and final disposal methods. Without reliable
records, environmental officers and administrators cannot accurately evaluate
recycling performance or resource recovery.

\subsection{Unsafe Handling and Disposal Practices}
Informal dismantling and disposal may expose workers and communities to
hazardous substances. Unsafe practices may include open burning of wires and
components, manual breaking of devices without protective equipment,
uncontrolled disposal of batteries, dumping of circuit boards and hazardous parts,
mixing electronic waste with municipal waste, and disposal near residential areas
or water sources. These practices may contribute to soil, air, and water pollution
and may create health risks for workers and surrounding communities.

\subsection{Limited Environmental Monitoring}
Environmental officers have limited access to real-time or centralized information
about electronic waste generation, collection, processing, and disposal.
Monitoring often depends on physical inspections and manually submitted
reports. This makes it difficult to identify electronic waste hotspots, uncollected
waste, delayed pickups, unauthorized processing activities, unsafe disposal
locations, poorly performing collection areas, and trends in waste generation and
recovery. The absence of data-driven monitoring reduces the effectiveness of
environmental decision-making.

\subsection{Inadequate Reporting and Analytics}
The existing process does not provide integrated dashboards and automated
reports. Administrators may find it difficult to generate accurate information on
registered users, electronic waste categories, monthly or annual waste volumes,
pickup-request performance, collection coverage, repaired and refurbished
devices, recycled waste and recovered materials, unresolved cases, and
environmental impact indicators. The absence of analytics limits strategic
planning, resource allocation, policy formulation, and performance evaluation.

\subsection{Limited User Accessibility}
The existing process does not provide consistent access through mobile, web, and
desktop platforms. Stakeholders with different operational responsibilities may
require different interfaces. For example, households may prefer mobile access,
administrators may require web dashboards, and operational staff may use
desktop systems. The absence of a cross-platform system restricts accessibility
and reduces stakeholder participation.

\subsection{Weak Security and Access Control}
Manual records and informal communication provide limited protection for
personal, operational, and environmental data. The existing process may not
provide secure authentication, user-role verification, permission-based access,
session management, audit trails, secure password recovery, or controlled access
to sensitive records. This weakness increases the risk of unauthorized access,
data exposure, record manipulation, and misuse of administrative functions.

\subsection{Limited Scalability}
Manual and fragmented processes are difficult to expand as the number of users,
collection requests, electronic waste records, facilities, and geographical locations
increases. The existing system does not provide cloud-based infrastructure
capable of supporting real-time synchronization, centralized data access,
automated deployment, and multi-platform growth.

\subsection{Limited Support for Circular Economy Activities}
The existing process often emphasizes disposal or material recovery without fully
coordinating reuse, repair, refurbishment, resale, and recycling. The lack of an
integrated circular economy platform reduces the ability to extend device
lifecycles, redistribute usable products, manage refurbished items, and recover
maximum economic value from electronic waste.

\subsection{Summary of Identified Problems}
The major weaknesses of the existing system can be summarized as follows:
(1)~absence of a centralized electronic waste management platform;
(2)~dependence on manual and fragmented records; (3)~limited lifecycle
traceability; (4)~inconsistent electronic waste classification; (5)~inefficient
pickup scheduling; (6)~lack of route optimization; (7)~weak coordination among
stakeholders; (8)~inadequate notification mechanisms; (9)~poor repair and
refurbishment documentation; (10)~incomplete recycling and resource-recovery
records; (11)~unsafe handling and disposal practices; (12)~limited environmental
monitoring; (13)~inadequate reporting and analytics; (14)~limited cross-platform
accessibility; (15)~weak security and access control; (16)~poor scalability; and
(17)~insufficient support for circular economy workflows.

These identified problems provide the basis for the analysis of the proposed
EcoTrace system presented in the next section. Figure~\ref{fig:existing-problems}
summarizes the major operational, technical, environmental, and administrative
weaknesses of the existing electronic waste management process.

\begin{figure}[H]
  \centering
  \imageplaceholder{Problems of the Existing Electronic Waste Management System}
  \caption{Problems of the Existing Electronic Waste Management System. Source: Author's analysis.}
  \label{fig:existing-problems}
\end{figure}

\section{Analysis of the Proposed System}

The proposed system is the EcoTrace: Intelligent Electronic Waste Recovery and
Circular Economy Management System. It is designed to address the
operational, technical, environmental, and administrative problems identified in
the existing electronic waste management process.

EcoTrace provides a centralized cross-platform environment for managing
electronic waste from the point of registration to its final outcome. The system
supports Android mobile applications, web applications, and Windows desktop
applications through Flutter. A FastAPI backend deployed on Render coordinates the
system's business logic and communication with external cloud services.

Cloud Firestore stores structured application data, Firebase Authentication
manages user identities, Firebase Cloud Messaging delivers notifications,
Cloudinary hosts electronic waste images, Google Maps Platform provides
location and route services, and PyTorch supports artificial-intelligence-based
image classification.

The proposed system supports households, businesses, collectors, drivers,
technicians, recyclers, environmental officers, administrators, and super
administrators. Each user category is assigned specific responsibilities and access
permissions.

\subsection{Objectives of the Proposed System}
The proposed system is designed to achieve the following objectives:

\begin{enumerate}[leftmargin=1.4cm]
  \item provide a centralized platform for electronic waste registration, collection,
    transportation, repair, refurbishment, recycling, and reporting
  \item improve electronic waste traceability throughout its lifecycle
  \item support AI-assisted classification of electronic waste images
  \item provide structured pickup scheduling and collector assignment
  \item improve location management and route planning using Google Maps
    Platform
  \item support secure authentication and role-based access control
  \item improve coordination among electronic waste stakeholders
  \item provide real-time notifications and status updates
  \item maintain repair, refurbishment, and recycling records
  \item support environmental monitoring and data-driven decision-making
  \item provide reports and analytics for operational and administrative use
  \item promote repair, reuse, refurbishment, recycling, and resource recovery in
    accordance with circular economy principles
\end{enumerate}

\subsection{Major Users of the Proposed System}
EcoTrace supports several user categories. Each user interacts with the system
according to assigned responsibilities.

\textbf{Guest User.} A guest user is an unauthenticated person who can access
public information and initiate account registration. The guest user may view
general information about EcoTrace, read electronic waste awareness content,
view publicly available collection-centre information, register a new account, and
access the login and password-recovery interfaces.

\textbf{Household User.} A household user represents an individual or family
generating electronic waste. The household user may create and manage a
personal profile, register electronic waste items, upload electronic waste images,
request AI-based waste classification, schedule pickup requests, select or record
pickup locations, monitor collection status, receive notifications, view repair,
refurbishment, and recycling outcomes, browse refurbished products, and
generate personal waste-management records.

\textbf{Business User.} A business user represents an organization, institution,
enterprise, or office that generates electronic waste. The business user may
maintain an organizational profile, register multiple electronic waste items,
upload asset images and supporting information, submit bulk pickup requests,
track organizational waste collections, monitor recycling and disposal outcomes,
download reports and certificates where available, and maintain disposal records
for accountability and compliance.

\textbf{Collector.} A collector is responsible for receiving, accepting, and
completing electronic waste pickup assignments. The collector may view assigned
pickup requests, review waste details and pickup locations, accept or reject
assignments, update collection status, confirm item collection, upload collection
evidence, communicate delays or exceptions, and transfer collected items to
drivers, centres, technicians, or recyclers.

\textbf{Driver.} A driver transports electronic waste between pickup locations,
collection centres, technicians, recyclers, and other processing facilities. The
driver may view assigned transportation tasks, access pickup and destination
locations, view suggested routes, update departure, transit, and delivery status,
record vehicle and delivery information, report transportation incidents, and
confirm delivery to the receiving facility.

\textbf{Technician.} A technician examines, diagnoses, repairs, and refurbishes
recoverable electronic devices. The technician may view assigned electronic
waste items, record inspection and diagnostic results, determine whether an item
is repairable, record required components and repair costs, update repair
progress, complete repair records, perform refurbishment assessments, record
quality-assurance results, and recommend recycling where repair is not practical.

\textbf{Recycler.} A recycler manages the dismantling, processing, material
recovery, and safe disposal of non-repairable electronic waste. The recycler may
receive electronic waste items for recycling, record waste category, quantity, and
condition, document dismantling activities, record recovered materials and
components, identify hazardous materials, record residual waste, update
recycling progress, confirm recycling completion, and generate recycling and
recovery reports.

\textbf{Environmental Officer.} An environmental officer monitors electronic
waste activities and environmental performance. The environmental officer may
view electronic waste locations and collection coverage, monitor collection,
transportation, repair, and recycling activities, inspect environmental reports,
identify delayed or unresolved cases, monitor recycling outcomes and recovered
materials, review potentially unsafe activities, analyse operational trends,
generate environmental reports, and support compliance monitoring and
decision-making.

\textbf{Administrator.} An administrator manages operational users, system
data, assignments, content, reports, and general system activities. The
administrator may manage user accounts and roles, verify organizations and
operational personnel, manage electronic waste categories, review pickup
requests, assign collectors and drivers, manage collection centres and facilities,
monitor repair, refurbishment, and recycling operations, manage marketplace
listings, send announcements and notifications, generate system reports, and
review audit and operational records.

\textbf{Super Administrator.} The super administrator has the highest level of
system authority and is responsible for security, configuration, high-level
administration, and oversight. The super administrator may create and manage
administrator accounts, configure global system settings, manage roles and
permissions, monitor security activities, review audit logs, manage API and
cloud-service configurations, suspend or restore critical accounts, manage
system-wide categories and policies, review system performance, and control
high-risk administrative operations.

\subsection{Major Modules of the Proposed System}
The proposed EcoTrace system consists of several interconnected modules.

\textbf{Authentication and User Management Module.} This module manages
account registration, login, logout, password recovery, email verification, user
profiles, roles, permissions, and account status. Firebase Authentication provides
identity management, while role and profile information is stored in Cloud
Firestore.

\textbf{Electronic Waste Registration Module.} This module enables households
and businesses to register electronic waste items. The user provides information
such as waste category, device name, brand and model where available,
condition, quantity, description, image, location, and preferred recovery option.
A unique record is created for each registered item to support lifecycle tracking.

\textbf{AI Classification Module.} This module analyses uploaded electronic
waste images using a PyTorch model. The image is transmitted to the FastAPI
backend, validated, preprocessed, and passed to the trained model. The system returns a
predicted waste category and confidence score. The final AI architecture remains
subject to model selection, training, testing, and performance evaluation.

\textbf{Pickup Scheduling Module.} This module enables users to request
collection services. The module records the electronic waste items included in
the request, preferred pickup date and time, pickup address and coordinates,
contact information, collection instructions, assigned collector or driver, and
request status.

\textbf{Collection and Assignment Module.} This module enables administrators
to review pickup requests and assign collectors and drivers. Collectors can accept
tasks, confirm collection, upload evidence, and update statuses. Administrators
can monitor pending, assigned, active, completed, delayed, and cancelled
requests.

\textbf{Transportation and Route Management Module.} This module supports
the transportation of electronic waste between users, collection centres,
technicians, recyclers, and processing facilities. Google Maps Platform provides
location display, distance calculation, route visualization, and route optimization.
Drivers can access assigned routes and update transportation status.

\textbf{Repair Management Module.} This module supports inspection,
diagnosis, repair planning, component replacement, cost recording,
repair-status updates, and final repair outcomes. Items that cannot be repaired
may be transferred to refurbishment or recycling workflows.

\textbf{Refurbishment Management Module.} This module supports the
restoration of recoverable devices to usable condition. It records refurbishment
activities, replaced components, testing results, quality grade, final condition,
availability for reuse or resale, and responsible technician.

\textbf{Recycling Management Module.} This module records electronic waste
received for recycling, dismantling activities, recovered materials, hazardous
components, residual waste, processing status, and final recycling outcomes. It
supports traceability and environmental reporting.

\textbf{Marketplace Module.} This module provides a controlled environment
for displaying refurbished devices, reusable components, and recovered
materials. Listings may contain product information, condition and quality
grade, images, availability, price, seller or facility details, and inspection or
refurbishment records.

\textbf{Notification Module.} This module delivers real-time alerts using
Firebase Cloud Messaging. Notifications may be generated for account
verification, pickup submission, request approval or rejection, collector
assignment, pickup reminders, collection completion, repair and refurbishment
updates, recycling completion, marketplace activity, environmental campaigns,
and system announcements.

\textbf{Reporting and Analytics Module.} This module generates operational,
administrative, and environmental reports. Reports may cover electronic waste
categories and quantities, pickup-request performance, collection coverage,
transportation activities, repair and refurbishment outcomes, recycling and
material recovery, marketplace activities, user activities, delayed or unresolved
cases, and environmental performance indicators.

\textbf{Environmental Monitoring Module.} This module provides dashboards
and map-based information for environmental officers. It supports the
monitoring of waste locations, collection coverage, facility activity, recycling
performance, recovery trends, and potential environmental concerns.

\textbf{Administration and System Configuration Module.} This module
supports user management, role configuration, waste categories, collection
centres, facilities, assignments, announcements, security controls, audit records,
integrations, and global system settings.

\subsection{Proposed System Workflow}
The general operational workflow of EcoTrace is as follows:

\begin{enumerate}[leftmargin=1.4cm]
  \item A household or business user creates an account and signs in.
  \item The user registers an electronic waste item and uploads an image.
  \item The AI module may classify the item and return a predicted category.
  \item The user submits a pickup request and records the pickup location.
  \item The administrator reviews the request and assigns a collector or driver.
  \item The collector receives the assignment and collects the electronic waste.
  \item The driver transports the item to an appropriate facility.
  \item The item is inspected and routed to repair, refurbishment, or recycling.
  \item Repairable devices are repaired or refurbished.
  \item Non-repairable items are dismantled and recycled.
  \item Recovered materials and reusable devices may be listed through the
    marketplace.
  \item Notifications are delivered during important workflow changes.
  \item Cloud Firestore maintains the lifecycle record.
  \item Administrators and environmental officers review reports, analytics, and
    operational activities.
\end{enumerate}

Figure~\ref{fig:proposed-workflow} illustrates the integrated electronic waste
management workflow proposed for EcoTrace.

\begin{figure}[H]
  \centering
  \imageplaceholder{Proposed EcoTrace Electronic Waste Management Workflow}
  \caption{Proposed EcoTrace Electronic Waste Management Workflow. Source: Author's design.}
  \label{fig:proposed-workflow}
\end{figure}

\section{Functional Requirements}
\label{sec:functional-requirements}

Functional requirements describe the services, operations, and behaviours that
the EcoTrace system shall provide. They define how the system shall respond to
user actions, process data, enforce workflows, communicate with external
services, and support the electronic waste management lifecycle.

The functional requirements are grouped according to the major modules and
user responsibilities of the system.

\subsection{Authentication and Account Management Requirements}
\begin{longtable}{@{}p{2.4cm}p{11cm}@{}}
\caption{Authentication and Account Management Requirements} \label{tab:fr-auth} \\
\toprule \textbf{ID} & \textbf{Requirement} \\ \midrule
\endfirsthead
\toprule \textbf{ID} & \textbf{Requirement} \\ \midrule
\endhead
FR-AUTH-1 & The system shall allow a guest user to create an account using a valid email address and password. \\
FR-AUTH-2 & The system shall validate required registration information before creating a user account. \\
FR-AUTH-3 & The system shall prevent the registration of duplicate email addresses. \\
FR-AUTH-4 & The system shall authenticate registered users using Firebase Authentication. \\
FR-AUTH-5 & The system shall reject invalid login credentials. \\
FR-AUTH-6 & The system shall allow authenticated users to log out securely. \\
FR-AUTH-7 & The system shall allow users to request password-reset instructions. \\
FR-AUTH-8 & The system shall support email-verification procedures where enabled. \\
FR-AUTH-9 & The system shall maintain an authenticated session until the user logs out or the session becomes invalid. \\
FR-AUTH-10 & The system shall redirect authenticated users to dashboards appropriate to their assigned roles. \\
FR-AUTH-11 & The system shall allow users to update permitted profile information. \\
FR-AUTH-12 & The system shall allow users to upload or update profile images where supported. \\
FR-AUTH-13 & The system shall allow administrators to activate, suspend, or restore user accounts. \\
FR-AUTH-14 & The system shall record the creation date and current status of each user account. \\
\bottomrule
\end{longtable}

\subsection{Role-Based Access Control Requirements}
\begin{longtable}{@{}p{2.4cm}p{11cm}@{}}
\caption{Role-Based Access Control Requirements} \label{tab:fr-rbac} \\
\toprule \textbf{ID} & \textbf{Requirement} \\ \midrule
\endfirsthead
\toprule \textbf{ID} & \textbf{Requirement} \\ \midrule
\endhead
FR-RBAC-1 & The system shall assign one or more approved roles to each registered user. \\
FR-RBAC-2 & The system shall restrict access to functions according to the authenticated user's role. \\
FR-RBAC-3 & The system shall prevent household users from accessing administrative functions. \\
FR-RBAC-4 & The system shall prevent collectors from accessing technician-only or recycler-only operations unless explicitly authorized. \\
FR-RBAC-5 & The system shall allow administrators to manage standard user roles. \\
FR-RBAC-6 & The system shall allow the super administrator to manage administrator accounts and high-level permissions. \\
FR-RBAC-7 & The system shall verify the user's role before processing protected backend requests. \\
FR-RBAC-8 & The system shall deny unauthorized operations and return an appropriate error response. \\
\bottomrule
\end{longtable}

\subsection{Electronic Waste Registration Requirements}
\begin{longtable}{@{}p{2.4cm}p{11cm}@{}}
\caption{Electronic Waste Registration Requirements} \label{tab:fr-ew} \\
\toprule \textbf{ID} & \textbf{Requirement} \\ \midrule
\endfirsthead
\toprule \textbf{ID} & \textbf{Requirement} \\ \midrule
\endhead
FR-EW-1 & The system shall allow household and business users to register electronic waste items. \\
FR-EW-2 & The system shall require the user to provide the electronic waste category, condition, quantity, and description. \\
FR-EW-3 & The system shall allow users to provide the brand and model where available. \\
FR-EW-4 & The system shall allow users to upload one or more electronic waste images where supported. \\
FR-EW-5 & The system shall upload electronic waste images through the FastAPI backend to Cloudinary. \\
FR-EW-6 & The system shall store the returned Cloudinary image URL with the corresponding electronic waste record. \\
FR-EW-7 & The system shall generate a unique identifier for each registered electronic waste item. \\
FR-EW-8 & The system shall record the user who registered the item. \\
FR-EW-9 & The system shall record the date and time of registration. \\
FR-EW-10 & The system shall record the current lifecycle status of the item. \\
FR-EW-11 & The system shall allow authorized users to view registered electronic waste details. \\
FR-EW-12 & The system shall allow the owner to update editable details before the item enters a restricted processing stage. \\
FR-EW-13 & The system shall prevent unauthorized users from modifying another user's electronic waste record. \\
FR-EW-14 & The system shall allow administrators to manage electronic waste categories. \\
\bottomrule
\end{longtable}

\subsection{Artificial Intelligence Classification Requirements}
\begin{longtable}{@{}p{2.4cm}p{11cm}@{}}
\caption{Artificial Intelligence Classification Requirements} \label{tab:fr-ai} \\
\toprule \textbf{ID} & \textbf{Requirement} \\ \midrule
\endfirsthead
\toprule \textbf{ID} & \textbf{Requirement} \\ \midrule
\endhead
FR-AI-1 & The system shall allow a user to submit an electronic waste image for classification. \\
FR-AI-2 & The FastAPI backend shall validate the submitted image format and size before processing. \\
FR-AI-3 & The system shall reject invalid, unsupported, or corrupted image files. \\
FR-AI-4 & The system shall preprocess the image according to the requirements of the selected PyTorch model. \\
FR-AI-5 & The system shall pass the preprocessed image to the trained PyTorch classification model. \\
FR-AI-6 & The system shall return a predicted electronic waste category. \\
FR-AI-7 & The system shall return a confidence score with the prediction where supported by the selected model. \\
FR-AI-8 & The system shall store the classification result with the relevant electronic waste record where the user accepts the result. \\
FR-AI-9 & The system shall allow the user or authorized staff to correct an incorrect classification. \\
FR-AI-10 & The system shall record whether a classification was generated by AI or entered manually. \\
FR-AI-11 & The system shall notify the user when asynchronous classification processing is completed. \\
FR-AI-12 & The system shall record classification errors for administrative review. \\
\bottomrule
\end{longtable}

\subsection{Pickup Request Requirements}
\begin{longtable}{@{}p{2.4cm}p{11cm}@{}}
\caption{Pickup Request Requirements} \label{tab:fr-pick} \\
\toprule \textbf{ID} & \textbf{Requirement} \\ \midrule
\endfirsthead
\toprule \textbf{ID} & \textbf{Requirement} \\ \midrule
\endhead
FR-PICK-1 & The system shall allow household and business users to create pickup requests. \\
FR-PICK-2 & The system shall allow one or more eligible electronic waste items to be included in a pickup request. \\
FR-PICK-3 & The system shall require a pickup address or geographic location. \\
FR-PICK-4 & The system shall allow the user to select a preferred pickup date and time. \\
FR-PICK-5 & The system shall allow the user to provide collection instructions. \\
FR-PICK-6 & The system shall record the requester's contact information. \\
FR-PICK-7 & The system shall generate a unique pickup-request identifier. \\
FR-PICK-8 & The system shall assign an initial status to each newly submitted request. \\
FR-PICK-9 & The system shall allow users to view the status of their pickup requests. \\
FR-PICK-10 & The system shall allow users to cancel eligible pickup requests before collection begins. \\
FR-PICK-11 & The system shall record the reason for cancellation where required. \\
FR-PICK-12 & The system shall prevent users from submitting invalid or incomplete pickup requests. \\
FR-PICK-13 & The system shall notify administrators when a new pickup request is submitted. \\
FR-PICK-14 & The system shall notify users when a pickup request is approved, rejected, scheduled, assigned, delayed, cancelled, or completed. \\
\bottomrule
\end{longtable}

\subsection{Collector and Driver Assignment Requirements}
\begin{longtable}{@{}p{2.4cm}p{11cm}@{}}
\caption{Collector and Driver Assignment Requirements} \label{tab:fr-asgn} \\
\toprule \textbf{ID} & \textbf{Requirement} \\ \midrule
\endfirsthead
\toprule \textbf{ID} & \textbf{Requirement} \\ \midrule
\endhead
FR-ASGN-1 & The system shall allow administrators to view unassigned pickup requests. \\
FR-ASGN-2 & The system shall allow administrators to assign a collector to an approved pickup request. \\
FR-ASGN-3 & The system shall allow administrators to assign a driver where transportation is required. \\
FR-ASGN-4 & The system shall record the date and time of each assignment. \\
FR-ASGN-5 & The system shall notify the assigned collector or driver. \\
FR-ASGN-6 & The system shall allow collectors to accept or reject assigned tasks. \\
FR-ASGN-7 & The system shall record the reason for task rejection where required. \\
FR-ASGN-8 & The system shall allow administrators to reassign rejected, cancelled, or delayed tasks. \\
FR-ASGN-9 & The system shall prevent the assignment of inactive or unauthorized operational users. \\
FR-ASGN-10 & The system shall display assignment history to authorized users. \\
\bottomrule
\end{longtable}

\subsection{Collection Management Requirements}
\begin{longtable}{@{}p{2.4cm}p{11cm}@{}}
\caption{Collection Management Requirements} \label{tab:fr-coll} \\
\toprule \textbf{ID} & \textbf{Requirement} \\ \midrule
\endfirsthead
\toprule \textbf{ID} & \textbf{Requirement} \\ \midrule
\endhead
FR-COLL-1 & The system shall allow collectors to view assigned pickup tasks. \\
FR-COLL-2 & The system shall display the waste details, requester information, pickup location, and schedule for each task. \\
FR-COLL-3 & The system shall allow collectors to update task status. \\
FR-COLL-4 & The system shall support statuses such as assigned, accepted, en route, arrived, collected, delayed, failed, and completed. \\
FR-COLL-5 & The system shall allow collectors to record the actual collection date and time. \\
FR-COLL-6 & The system shall allow collectors to confirm the quantity and condition of collected items. \\
FR-COLL-7 & The system shall allow collectors to upload collection evidence where required. \\
FR-COLL-8 & The system shall allow collectors to report missing, unavailable, or incorrectly described waste items. \\
FR-COLL-9 & The system shall update the lifecycle status of collected electronic waste items. \\
FR-COLL-10 & The system shall notify the requester after successful collection. \\
\bottomrule
\end{longtable}

\subsection{Transportation and Route Management Requirements}
\begin{longtable}{@{}p{2.4cm}p{11cm}@{}}
\caption{Transportation and Route Management Requirements} \label{tab:fr-route} \\
\toprule \textbf{ID} & \textbf{Requirement} \\ \midrule
\endfirsthead
\toprule \textbf{ID} & \textbf{Requirement} \\ \midrule
\endhead
FR-ROUTE-1 & The system shall capture or store geographic coordinates for pickup and destination locations. \\
FR-ROUTE-2 & The system shall display pickup locations using Google Maps Platform. \\
FR-ROUTE-3 & The system shall allow users to search for relevant collection centres where supported. \\
FR-ROUTE-4 & The system shall calculate route distance and estimated travel time where supported by the selected Google Maps service. \\
FR-ROUTE-5 & The system shall display suggested routes to authorized collectors and drivers. \\
FR-ROUTE-6 & The system shall support route optimization for multiple collection locations where enabled. \\
FR-ROUTE-7 & The system shall allow drivers to view assigned transportation tasks. \\
FR-ROUTE-8 & The system shall allow drivers to update departure, transit, arrival, and delivery statuses. \\
FR-ROUTE-9 & The system shall record the origin and destination of each transportation task. \\
FR-ROUTE-10 & The system shall allow drivers to report delays, incidents, or route exceptions. \\
FR-ROUTE-11 & The system shall record delivery confirmation at the receiving facility. \\
FR-ROUTE-12 & The system shall update the electronic waste lifecycle record after delivery. \\
\bottomrule
\end{longtable}

\subsection{Repair Management Requirements}
\begin{longtable}{@{}p{2.4cm}p{11cm}@{}}
\caption{Repair Management Requirements} \label{tab:fr-rep} \\
\toprule \textbf{ID} & \textbf{Requirement} \\ \midrule
\endfirsthead
\toprule \textbf{ID} & \textbf{Requirement} \\ \midrule
\endhead
FR-REP-1 & The system shall allow authorized personnel to assign an item to a technician. \\
FR-REP-2 & The system shall allow technicians to view assigned items. \\
FR-REP-3 & The system shall allow technicians to record inspection and diagnostic findings. \\
FR-REP-4 & The system shall allow technicians to indicate whether an item is repairable. \\
FR-REP-5 & The system shall allow technicians to record required parts, estimated cost, and expected completion date. \\
FR-REP-6 & The system shall allow technicians to update repair status. \\
FR-REP-7 & The system shall allow technicians to record replaced components. \\
FR-REP-8 & The system shall allow technicians to record testing results after repair. \\
FR-REP-9 & The system shall allow technicians to mark a repair as successful, unsuccessful, cancelled, or transferred. \\
FR-REP-10 & The system shall transfer non-repairable items to refurbishment or recycling workflows where appropriate. \\
FR-REP-11 & The system shall maintain the repair history of each item. \\
FR-REP-12 & The system shall notify relevant users when repair status changes. \\
\bottomrule
\end{longtable}

\subsection{Refurbishment Management Requirements}
\begin{longtable}{@{}p{2.4cm}p{11cm}@{}}
\caption{Refurbishment Management Requirements} \label{tab:fr-ref} \\
\toprule \textbf{ID} & \textbf{Requirement} \\ \midrule
\endfirsthead
\toprule \textbf{ID} & \textbf{Requirement} \\ \midrule
\endhead
FR-REF-1 & The system shall allow eligible items to enter the refurbishment workflow. \\
FR-REF-2 & The system shall allow technicians to record refurbishment activities. \\
FR-REF-3 & The system shall allow technicians to record replaced or upgraded components. \\
FR-REF-4 & The system shall allow technicians to record functional testing results. \\
FR-REF-5 & The system shall allow technicians to assign a quality grade. \\
FR-REF-6 & The system shall allow technicians to record the final condition of the refurbished item. \\
FR-REF-7 & The system shall allow authorized users to approve or reject a refurbished item for reuse or resale. \\
FR-REF-8 & The system shall store the refurbishment history of each item. \\
FR-REF-9 & The system shall allow approved refurbished items to be listed in the marketplace. \\
FR-REF-10 & The system shall notify relevant users when refurbishment is completed. \\
\bottomrule
\end{longtable}

\subsection{Recycling Management Requirements}
\begin{longtable}{@{}p{2.4cm}p{11cm}@{}}
\caption{Recycling Management Requirements} \label{tab:fr-rec} \\
\toprule \textbf{ID} & \textbf{Requirement} \\ \midrule
\endfirsthead
\toprule \textbf{ID} & \textbf{Requirement} \\ \midrule
\endhead
FR-REC-1 & The system shall allow authorized users to assign non-repairable items to a recycler. \\
FR-REC-2 & The system shall allow recyclers to view assigned items. \\
FR-REC-3 & The system shall allow recyclers to record the date and quantity of waste received. \\
FR-REC-4 & The system shall allow recyclers to record dismantling activities. \\
FR-REC-5 & The system shall allow recyclers to record recovered materials and components. \\
FR-REC-6 & The system shall allow recyclers to identify hazardous materials. \\
FR-REC-7 & The system shall allow recyclers to record residual waste. \\
FR-REC-8 & The system shall allow recyclers to record the final disposal method. \\
FR-REC-9 & The system shall allow recyclers to update recycling status. \\
FR-REC-10 & The system shall record the date and responsible personnel when recycling is completed. \\
FR-REC-11 & The system shall update the electronic waste lifecycle status after recycling. \\
FR-REC-12 & The system shall generate recycling records and reports. \\
FR-REC-13 & The system shall notify relevant users when recycling is completed. \\
\bottomrule
\end{longtable}

\subsection{Marketplace Requirements}
\begin{longtable}{@{}p{2.4cm}p{11cm}@{}}
\caption{Marketplace Requirements} \label{tab:fr-mkt} \\
\toprule \textbf{ID} & \textbf{Requirement} \\ \midrule
\endfirsthead
\toprule \textbf{ID} & \textbf{Requirement} \\ \midrule
\endhead
FR-MKT-1 & The system shall allow authorized users to create marketplace listings for approved refurbished devices, reusable parts, or recovered materials. \\
FR-MKT-2 & The system shall require a title, description, category, condition, price, availability, and image for each listing. \\
FR-MKT-3 & The system shall link a listing to its refurbishment or recovery record where applicable. \\
FR-MKT-4 & The system shall allow users to browse available listings. \\
FR-MKT-5 & The system shall allow users to search and filter listings. \\
FR-MKT-6 & The system shall display product condition and quality information. \\
FR-MKT-7 & The system shall allow administrators to review, approve, suspend, or remove listings. \\
FR-MKT-8 & The system shall prevent inactive or rejected listings from being displayed publicly. \\
FR-MKT-9 & The system shall record listing status and creation date. \\
FR-MKT-10 & The system shall support purchase or enquiry workflows according to the implemented marketplace scope. \\
\bottomrule
\end{longtable}

\subsection{Notification Requirements}
\begin{longtable}{@{}p{2.4cm}p{11cm}@{}}
\caption{Notification Requirements} \label{tab:fr-not} \\
\toprule \textbf{ID} & \textbf{Requirement} \\ \midrule
\endfirsthead
\toprule \textbf{ID} & \textbf{Requirement} \\ \midrule
\endhead
FR-NOT-1 & The system shall send notifications through Firebase Cloud Messaging where supported. \\
FR-NOT-2 & The system shall send notifications to specific users based on account identifiers or registered device tokens. \\
FR-NOT-3 & The system shall support notifications for pickup, collection, transportation, repair, refurbishment, recycling, marketplace, security, and administrative events. \\
FR-NOT-4 & The system shall store notification records in Cloud Firestore. \\
FR-NOT-5 & The system shall display unread and read notification status. \\
FR-NOT-6 & The system shall allow users to open a notification and view the related record where applicable. \\
FR-NOT-7 & The system shall allow authorized administrators to send system announcements. \\
FR-NOT-8 & The system shall prevent unauthorized users from sending global notifications. \\
FR-NOT-9 & The system shall record notification creation and delivery information where available. \\
\bottomrule
\end{longtable}

\subsection{Reporting and Analytics Requirements}
\begin{longtable}{@{}p{2.4cm}p{11cm}@{}}
\caption{Reporting and Analytics Requirements} \label{tab:fr-rpt} \\
\toprule \textbf{ID} & \textbf{Requirement} \\ \midrule
\endfirsthead
\toprule \textbf{ID} & \textbf{Requirement} \\ \midrule
\endhead
FR-RPT-1 & The system shall generate reports based on authorized user roles. \\
FR-RPT-2 & The system shall provide summary statistics for registered electronic waste. \\
FR-RPT-3 & The system shall provide statistics on pickup requests and collection performance. \\
FR-RPT-4 & The system shall provide reports on repair and refurbishment outcomes. \\
FR-RPT-5 & The system shall provide reports on recycling activities and recovered materials. \\
FR-RPT-6 & The system shall provide reports on users, roles, and operational activities. \\
FR-RPT-7 & The system shall allow authorized users to filter reports by date, status, category, location, or responsible personnel where applicable. \\
FR-RPT-8 & The system shall display analytical information using tables, charts, indicators, and maps where appropriate. \\
FR-RPT-9 & The system shall allow authorized users to export reports in PDF or other supported formats. \\
FR-RPT-10 & The system shall prevent unauthorized users from viewing restricted reports. \\
\bottomrule
\end{longtable}

\subsection{Environmental Monitoring Requirements}
\begin{longtable}{@{}p{2.4cm}p{11cm}@{}}
\caption{Environmental Monitoring Requirements} \label{tab:fr-env} \\
\toprule \textbf{ID} & \textbf{Requirement} \\ \midrule
\endfirsthead
\toprule \textbf{ID} & \textbf{Requirement} \\ \midrule
\endhead
FR-ENV-1 & The system shall provide environmental officers with a dedicated dashboard. \\
FR-ENV-2 & The system shall display electronic waste collection locations on a map. \\
FR-ENV-3 & The system shall allow environmental officers to monitor collection, repair, refurbishment, and recycling activities. \\
FR-ENV-4 & The system shall display unresolved, delayed, or potentially unsafe cases. \\
FR-ENV-5 & The system shall provide statistics on electronic waste categories and quantities. \\
FR-ENV-6 & The system shall provide statistics on recovered materials and recycling outcomes. \\
FR-ENV-7 & The system shall allow environmental officers to generate environmental reports. \\
FR-ENV-8 & The system shall allow environmental officers to review facility and operational records. \\
FR-ENV-9 & The system shall allow authorized officers to record observations, inspection notes, or compliance findings where implemented. \\
\bottomrule
\end{longtable}

\subsection{Administration Requirements}
\begin{longtable}{@{}p{2.4cm}p{11cm}@{}}
\caption{Administration Requirements} \label{tab:fr-adm} \\
\toprule \textbf{ID} & \textbf{Requirement} \\ \midrule
\endfirsthead
\toprule \textbf{ID} & \textbf{Requirement} \\ \midrule
\endhead
FR-ADM-1 & The system shall provide administrators with an operational dashboard. \\
FR-ADM-2 & The system shall allow administrators to manage users and account statuses. \\
FR-ADM-3 & The system shall allow administrators to manage user roles within their authority. \\
FR-ADM-4 & The system shall allow administrators to manage electronic waste categories. \\
FR-ADM-5 & The system shall allow administrators to manage collection centres, repair facilities, and recycling facilities. \\
FR-ADM-6 & The system shall allow administrators to review pickup requests. \\
FR-ADM-7 & The system shall allow administrators to assign collectors and drivers. \\
FR-ADM-8 & The system shall allow administrators to monitor active workflows. \\
FR-ADM-9 & The system shall allow administrators to manage marketplace listings. \\
FR-ADM-10 & The system shall allow administrators to publish announcements. \\
FR-ADM-11 & The system shall allow administrators to generate operational reports. \\
FR-ADM-12 & The system shall record significant administrative actions. \\
\bottomrule
\end{longtable}

\subsection{Super Administration and Security Requirements}
\begin{longtable}{@{}p{2.4cm}p{11cm}@{}}
\caption{Super Administration and Security Requirements} \label{tab:fr-sadm} \\
\toprule \textbf{ID} & \textbf{Requirement} \\ \midrule
\endfirsthead
\toprule \textbf{ID} & \textbf{Requirement} \\ \midrule
\endhead
FR-SADM-1 & The system shall provide the super administrator with access to high-level configuration functions. \\
FR-SADM-2 & The system shall allow the super administrator to create, update, suspend, and restore administrator accounts. \\
FR-SADM-3 & The system shall allow the super administrator to manage roles and permissions. \\
FR-SADM-4 & The system shall allow the super administrator to review audit records. \\
FR-SADM-5 & The system shall allow the super administrator to configure system-wide settings. \\
FR-SADM-6 & The system shall allow the super administrator to manage external service configuration values through secure deployment mechanisms. \\
FR-SADM-7 & The system shall prevent standard administrators from performing restricted super-administrator operations. \\
FR-SADM-8 & The system shall record high-risk administrative actions. \\
FR-SADM-9 & The system shall allow the super administrator to review system health and integration status where implemented. \\
\bottomrule
\end{longtable}

\subsection{Audit and Activity Tracking Requirements}
\begin{longtable}{@{}p{2.4cm}p{11cm}@{}}
\caption{Audit and Activity Tracking Requirements} \label{tab:fr-aud} \\
\toprule \textbf{ID} & \textbf{Requirement} \\ \midrule
\endfirsthead
\toprule \textbf{ID} & \textbf{Requirement} \\ \midrule
\endhead
FR-AUD-1 & The system shall record significant user and administrative activities. \\
FR-AUD-2 & The system shall record the actor, action, affected record, and timestamp where applicable. \\
FR-AUD-3 & The system shall maintain status histories for electronic waste, pickup, transportation, repair, refurbishment, and recycling records. \\
FR-AUD-4 & The system shall restrict access to audit information. \\
FR-AUD-5 & The system shall prevent standard users from modifying audit records. \\
FR-AUD-6 & The system shall allow authorized administrators to search or filter audit records. \\
\bottomrule
\end{longtable}

\subsection{External Service Integration Requirements}
\begin{longtable}{@{}p{2.4cm}p{11cm}@{}}
\caption{External Service Integration Requirements} \label{tab:fr-int} \\
\toprule \textbf{ID} & \textbf{Requirement} \\ \midrule
\endfirsthead
\toprule \textbf{ID} & \textbf{Requirement} \\ \midrule
\endhead
FR-INT-1 & The system shall communicate with the FastAPI backend using secure HTTP requests. \\
FR-INT-2 & The system shall exchange structured data using JSON. \\
FR-INT-3 & The system shall integrate with Firebase Authentication for user identity management. \\
FR-INT-4 & The system shall integrate with Cloud Firestore for structured application data. \\
FR-INT-5 & The system shall integrate with Cloudinary for electronic waste image hosting. \\
FR-INT-6 & The system shall integrate with Firebase Cloud Messaging for push notifications. \\
FR-INT-7 & The system shall integrate with Google Maps Platform for mapping and routing services. \\
FR-INT-8 & The system shall integrate with PyTorch for AI-based image classification. \\
FR-INT-9 & The FastAPI backend shall be deployable on the Render cloud platform. \\
FR-INT-10 & The system shall handle unavailable or failed external-service responses without exposing sensitive technical details to standard users. \\
\bottomrule
\end{longtable}

\subsection{Functional Requirements Summary}
The functional requirements define the operations that EcoTrace shall perform
across authentication, waste registration, artificial intelligence, pickup
scheduling, collection, transportation, repair, refurbishment, recycling,
marketplace management, notifications, reporting, environmental monitoring,
administration, security, audit tracking, and external service integration.

These requirements provide the basis for the system models, database design,
interface design, implementation, and test cases described in the remaining
sections of this chapter and in Chapter~\ref{ch:implementation}.

\section{Non-Functional Requirements}
\label{sec:nfr}

Non-functional requirements describe the quality attributes, constraints, and
operational standards that the EcoTrace system shall satisfy. Unlike functional
requirements, which define what the system does, non-functional requirements
specify how effectively, securely, reliably, and efficiently the system shall perform
its functions.

The non-functional requirements of EcoTrace cover performance, availability,
reliability, security, privacy, usability, accessibility, scalability, maintainability,
compatibility, data integrity, backup, recovery, auditability, portability, and
external-service resilience.

\subsection{Performance Requirements}
\begin{longtable}{@{}p{2.6cm}p{10.8cm}@{}}
\caption{Performance Requirements} \label{tab:nfr-perf} \\
\toprule \textbf{ID} & \textbf{Requirement} \\ \midrule
\endfirsthead
\toprule \textbf{ID} & \textbf{Requirement} \\ \midrule
\endhead
NFR-PERF-1 & The system shall provide acceptable response times for normal user operations under the expected academic demonstration workload. \\
NFR-PERF-2 & Standard interface actions such as navigation, record retrieval, and form submission should complete without unnecessary delay when network connectivity and external services are available. \\
NFR-PERF-3 & The FastAPI backend shall process valid API requests efficiently and return structured responses. \\
NFR-PERF-4 & The system shall avoid blocking the user interface during image upload, notification processing, route calculation, report generation, or AI inference. \\
NFR-PERF-5 & Long-running operations shall display progress indicators, loading states, or processing messages. \\
NFR-PERF-6 & The system shall minimize unnecessary network requests and repeated database reads. \\
NFR-PERF-7 & Images shall be compressed or optimized before or during cloud delivery where appropriate. \\
NFR-PERF-8 & Database queries shall use suitable filters and indexes to reduce unnecessary data retrieval. \\
NFR-PERF-9 & The AI classification service shall return a result within a reasonable period based on the selected model and available Render compute resources. \\
NFR-PERF-10 & Performance measurements such as API response time, page-loading time, image-upload time, and AI inference time shall be evaluated during system testing. \\
\bottomrule
\end{longtable}

\subsection{Availability Requirements}
\begin{longtable}{@{}p{2.6cm}p{10.8cm}@{}}
\caption{Availability Requirements} \label{tab:nfr-avl} \\
\toprule \textbf{ID} & \textbf{Requirement} \\ \midrule
\endfirsthead
\toprule \textbf{ID} & \textbf{Requirement} \\ \midrule
\endhead
NFR-AVL-1 & The system shall remain accessible whenever the Flutter applications, backend service, Internet connection, and required third-party services are available. \\
NFR-AVL-2 & The system shall provide informative feedback when the backend or an external service is temporarily unavailable. \\
NFR-AVL-3 & The system shall not display a successful operation when the corresponding backend request has failed. \\
NFR-AVL-4 & The system shall support recovery after temporary network disconnection. \\
NFR-AVL-5 & The system shall preserve valid unsent form data where practical when a temporary connection failure occurs. \\
NFR-AVL-6 & The deployed FastAPI service shall expose a health-check endpoint where implemented. \\
NFR-AVL-7 & The system shall recognize that the Render Free Web Service may enter an inactive state and produce an initial cold-start delay. \\
NFR-AVL-8 & Production deployment should use hosting resources appropriate to the expected availability requirements. \\
\bottomrule
\end{longtable}

\subsection{Reliability Requirements}
\begin{longtable}{@{}p{2.6cm}p{10.8cm}@{}}
\caption{Reliability Requirements} \label{tab:nfr-rel} \\
\toprule \textbf{ID} & \textbf{Requirement} \\ \midrule
\endfirsthead
\toprule \textbf{ID} & \textbf{Requirement} \\ \midrule
\endhead
NFR-REL-1 & The system shall process valid operations consistently and produce predictable results. \\
NFR-REL-2 & The system shall validate input before saving or processing data. \\
NFR-REL-3 & The system shall prevent incomplete records from being treated as completed operations. \\
NFR-REL-4 & The system shall preserve the consistency of related electronic waste, pickup, assignment, transportation, repair, refurbishment, and recycling records. \\
NFR-REL-5 & The system shall handle failed requests without corrupting existing data. \\
NFR-REL-6 & The system shall display clear error messages when an operation cannot be completed. \\
NFR-REL-7 & The system shall avoid duplicate record creation when a user submits the same operation repeatedly due to network delay. \\
NFR-REL-8 & Critical status transitions shall be validated before being saved. \\
NFR-REL-9 & The system shall record failed integration operations where administrative review is required. \\
NFR-REL-10 & Reliability shall be assessed through repeated functional, integration, and system tests. \\
\bottomrule
\end{longtable}

\subsection{Security Requirements}
\begin{longtable}{@{}p{2.6cm}p{10.8cm}@{}}
\caption{Security Requirements} \label{tab:nfr-sec} \\
\toprule \textbf{ID} & \textbf{Requirement} \\ \midrule
\endfirsthead
\toprule \textbf{ID} & \textbf{Requirement} \\ \midrule
\endhead
NFR-SEC-1 & The system shall require authentication before granting access to protected functions. \\
NFR-SEC-2 & The system shall enforce role-based access control for all restricted operations. \\
NFR-SEC-3 & The backend shall verify authorization independently of the client interface. \\
NFR-SEC-4 & The system shall transmit sensitive data through HTTPS. \\
NFR-SEC-5 & Passwords shall be managed through Firebase Authentication rather than stored directly in Cloud Firestore. \\
NFR-SEC-6 & The system shall not expose passwords, private keys, service credentials, API secrets, or authentication tokens in source code, client logs, or error messages. \\
NFR-SEC-7 & Backend credentials and external-service secrets shall be stored using environment variables or equivalent secure configuration mechanisms. \\
NFR-SEC-8 & Google Maps API keys shall be restricted according to the authorized application and enabled APIs. \\
NFR-SEC-9 & The system shall validate and sanitize data received by backend endpoints. \\
NFR-SEC-10 & Uploaded files shall be checked for allowed format, size, and content type before processing. \\
NFR-SEC-11 & The system shall prevent unauthorized modification of another user's records. \\
NFR-SEC-12 & Administrative and super-administrative functions shall require appropriate privileges. \\
NFR-SEC-13 & The system shall record significant security-sensitive actions. \\
NFR-SEC-14 & User sessions shall become invalid when authentication tokens expire, are revoked, or the user logs out. \\
NFR-SEC-15 & The system shall avoid exposing internal stack traces and technical configuration details to standard users. \\
NFR-SEC-16 & Firestore Security Rules shall restrict direct database access according to authentication status, ownership, and role. \\
\bottomrule
\end{longtable}

\subsection{Privacy Requirements}
\begin{longtable}{@{}p{2.6cm}p{10.8cm}@{}}
\caption{Privacy Requirements} \label{tab:nfr-priv} \\
\toprule \textbf{ID} & \textbf{Requirement} \\ \midrule
\endfirsthead
\toprule \textbf{ID} & \textbf{Requirement} \\ \midrule
\endhead
NFR-PRIV-1 & The system shall collect only personal and operational information required for electronic waste management activities. \\
NFR-PRIV-2 & Personal information shall be accessible only to authorized users. \\
NFR-PRIV-3 & Users shall not be able to view private profiles or records belonging to unrelated users. \\
NFR-PRIV-4 & Location information shall be used only for authorized pickup, transportation, mapping, and monitoring functions. \\
NFR-PRIV-5 & Uploaded electronic waste images shall be linked only to the relevant authorized records. \\
NFR-PRIV-6 & Reports intended for general viewing shall avoid exposing unnecessary personal information. \\
NFR-PRIV-7 & Administrative exports shall be restricted to authorized roles. \\
NFR-PRIV-8 & The system shall support account deactivation or deletion procedures according to the implemented policy. \\
NFR-PRIV-9 & Audit records shall not expose confidential information beyond what is required for accountability. \\
\bottomrule
\end{longtable}

\subsection{Usability Requirements}
\begin{longtable}{@{}p{2.6cm}p{10.8cm}@{}}
\caption{Usability Requirements} \label{tab:nfr-usa} \\
\toprule \textbf{ID} & \textbf{Requirement} \\ \midrule
\endfirsthead
\toprule \textbf{ID} & \textbf{Requirement} \\ \midrule
\endhead
NFR-USA-1 & The system shall provide a clear and consistent user interface across Android, Web, and Windows platforms. \\
NFR-USA-2 & Navigation labels, buttons, forms, icons, and status indicators shall use understandable terminology. \\
NFR-USA-3 & The system shall display dashboard content relevant to the user's role. \\
NFR-USA-4 & Forms shall identify required fields clearly. \\
NFR-USA-5 & Validation messages shall explain how the user can correct invalid input. \\
NFR-USA-6 & Destructive operations shall require confirmation where appropriate. \\
NFR-USA-7 & The system shall provide visible feedback after successful or failed operations. \\
NFR-USA-8 & Status values shall use consistent names and visual indicators throughout the system. \\
NFR-USA-9 & Common operations shall require a reasonable number of user steps. \\
NFR-USA-10 & The system shall maintain consistent layout, typography, spacing, and component behaviour. \\
NFR-USA-11 & User interfaces shall be designed for users with different levels of technical experience. \\
NFR-USA-12 & The system shall provide helpful empty-state messages when no data is available. \\
NFR-USA-13 & Loading indicators shall be displayed while data is being retrieved or processed. \\
NFR-USA-14 & Date, time, currency, quantity, and location information shall be presented in understandable formats. \\
\bottomrule
\end{longtable}

\subsection{Accessibility Requirements}
\begin{longtable}{@{}p{2.6cm}p{10.8cm}@{}}
\caption{Accessibility Requirements} \label{tab:nfr-acc} \\
\toprule \textbf{ID} & \textbf{Requirement} \\ \midrule
\endfirsthead
\toprule \textbf{ID} & \textbf{Requirement} \\ \midrule
\endhead
NFR-ACC-1 & The system should use readable font sizes and sufficient visual contrast. \\
NFR-ACC-2 & Important information shall not be communicated through colour alone. \\
NFR-ACC-3 & Interactive elements should provide clear labels or tooltips. \\
NFR-ACC-4 & Touch targets on mobile interfaces should be large enough for reliable interaction. \\
NFR-ACC-5 & Forms should follow a logical focus and reading order. \\
NFR-ACC-6 & Images that communicate essential information should have descriptive alternatives where supported. \\
NFR-ACC-7 & The Web application should support keyboard navigation for major operations where practical. \\
NFR-ACC-8 & Error and status messages shall be visible and understandable. \\
NFR-ACC-9 & Interfaces should adapt to different screen sizes and display scaling settings. \\
\bottomrule
\end{longtable}

\subsection{Scalability Requirements}
\begin{longtable}{@{}p{2.6cm}p{10.8cm}@{}}
\caption{Scalability Requirements} \label{tab:nfr-scal} \\
\toprule \textbf{ID} & \textbf{Requirement} \\ \midrule
\endfirsthead
\toprule \textbf{ID} & \textbf{Requirement} \\ \midrule
\endhead
NFR-SCAL-1 & The system architecture shall support growth in users, electronic waste records, pickup requests, images, notifications, and reports. \\
NFR-SCAL-2 & Cloud Firestore collections shall be structured to support future data expansion. \\
NFR-SCAL-3 & Backend services shall use modular components that can be scaled or deployed independently where necessary. \\
NFR-SCAL-4 & The system shall avoid design decisions that require all records to be loaded into memory simultaneously. \\
NFR-SCAL-5 & Reports and dashboards shall support filtered or paginated data retrieval where large datasets are expected. \\
NFR-SCAL-6 & Image files shall be stored in Cloudinary rather than in the Render local filesystem. \\
NFR-SCAL-7 & The architecture shall support migration from the Render Free tier to paid or dedicated infrastructure. \\
NFR-SCAL-8 & The AI service shall allow replacement or upgrading of the selected model without redesigning the complete client application. \\
NFR-SCAL-9 & New locations, collection centres, facilities, user roles, and electronic waste categories shall be addable without major architectural changes. \\
\bottomrule
\end{longtable}

\subsection{Maintainability Requirements}
\begin{longtable}{@{}p{2.6cm}p{10.8cm}@{}}
\caption{Maintainability Requirements} \label{tab:nfr-main} \\
\toprule \textbf{ID} & \textbf{Requirement} \\ \midrule
\endfirsthead
\toprule \textbf{ID} & \textbf{Requirement} \\ \midrule
\endhead
NFR-MAIN-1 & The system source code shall be organized into logical modules. \\
NFR-MAIN-2 & Business logic shall be separated from user-interface components where practical. \\
NFR-MAIN-3 & Backend routes, services, schemas, and integration components shall be structured clearly. \\
NFR-MAIN-4 & Repeated logic should be implemented through reusable functions or components. \\
NFR-MAIN-5 & Configuration values shall be separated from application logic. \\
NFR-MAIN-6 & Source code shall use clear and consistent naming conventions. \\
NFR-MAIN-7 & Complex functions and modules shall include useful documentation or comments. \\
NFR-MAIN-8 & The project shall use version control through Git and GitHub. \\
NFR-MAIN-9 & Changes shall be tested before deployment to the main environment. \\
NFR-MAIN-10 & Database fields, API responses, and status values shall follow consistent naming conventions. \\
NFR-MAIN-11 & External-service integrations shall be isolated sufficiently to allow replacement or modification. \\
NFR-MAIN-12 & The system shall support future correction, extension, and module replacement without requiring complete redevelopment. \\
\bottomrule
\end{longtable}

\subsection{Compatibility Requirements}
\begin{longtable}{@{}p{2.6cm}p{10.8cm}@{}}
\caption{Compatibility Requirements} \label{tab:nfr-comp} \\
\toprule \textbf{ID} & \textbf{Requirement} \\ \midrule
\endfirsthead
\toprule \textbf{ID} & \textbf{Requirement} \\ \midrule
\endhead
NFR-COMP-1 & The mobile application shall support the targeted Android versions defined during implementation. \\
NFR-COMP-2 & The Web application shall operate in current mainstream browsers used during testing. \\
NFR-COMP-3 & The Windows application shall operate on the targeted Windows versions defined during implementation. \\
NFR-COMP-4 & Client applications shall communicate with the FastAPI backend through standard HTTP or HTTPS requests. \\
NFR-COMP-5 & API data shall use JSON unless another documented format is required. \\
NFR-COMP-6 & Images shall use supported formats such as JPEG, JPG, PNG, or other approved formats. \\
NFR-COMP-7 & The system shall use platform-responsive layouts to reduce interface overflow and display errors. \\
NFR-COMP-8 & Platform-specific features shall provide suitable alternatives where they are unavailable on another platform. \\
NFR-COMP-9 & The system shall maintain consistent core workflows across Android, Web, and Windows. \\
\bottomrule
\end{longtable}

\subsection{Data Integrity Requirements}
\begin{longtable}{@{}p{2.6cm}p{10.8cm}@{}}
\caption{Data Integrity Requirements} \label{tab:nfr-data} \\
\toprule \textbf{ID} & \textbf{Requirement} \\ \midrule
\endfirsthead
\toprule \textbf{ID} & \textbf{Requirement} \\ \midrule
\endhead
NFR-DATA-1 & Each important record shall have a unique identifier. \\
NFR-DATA-2 & Required fields shall be validated before data is stored. \\
NFR-DATA-3 & Relationships between user, waste, pickup, assignment, transport, repair, refurbishment, recycling, and marketplace records shall be maintained correctly. \\
NFR-DATA-4 & Status transitions shall follow approved workflow rules. \\
NFR-DATA-5 & Timestamps shall be recorded for significant record creation and update activities. \\
NFR-DATA-6 & The system shall avoid orphaned records where practical. \\
NFR-DATA-7 & Deleted or deactivated records shall be handled according to the implemented retention policy. \\
NFR-DATA-8 & Audit and history records shall preserve evidence of important changes. \\
NFR-DATA-9 & The system shall not overwrite important historical information without authorization. \\
NFR-DATA-10 & Data retrieved by different platforms shall remain consistent after successful synchronization. \\
\bottomrule
\end{longtable}

\subsection{Backup and Recovery Requirements}
\begin{longtable}{@{}p{2.6cm}p{10.8cm}@{}}
\caption{Backup and Recovery Requirements} \label{tab:nfr-back} \\
\toprule \textbf{ID} & \textbf{Requirement} \\ \midrule
\endfirsthead
\toprule \textbf{ID} & \textbf{Requirement} \\ \midrule
\endhead
NFR-BACK-1 & Important system data shall be stored in managed cloud services rather than only on local client devices. \\
NFR-BACK-2 & Structured data shall be stored in Cloud Firestore. \\
NFR-BACK-3 & Uploaded images shall be stored in Cloudinary. \\
NFR-BACK-4 & The system shall not depend on the ephemeral Render filesystem for permanent data storage. \\
NFR-BACK-5 & Configuration and source code shall be maintained in version control. \\
NFR-BACK-6 & Production deployment should provide documented backup and restoration procedures. \\
NFR-BACK-7 & The system shall support recovery from temporary client or network failure without corrupting saved cloud data. \\
NFR-BACK-8 & Critical records should include timestamps and status histories to support reconstruction of operational events. \\
NFR-BACK-9 & Backup and restoration procedures shall be tested before large-scale production deployment. \\
\bottomrule
\end{longtable}

\subsection{Auditability Requirements}
\begin{longtable}{@{}p{2.6cm}p{10.8cm}@{}}
\caption{Auditability Requirements} \label{tab:nfr-aud} \\
\toprule \textbf{ID} & \textbf{Requirement} \\ \midrule
\endfirsthead
\toprule \textbf{ID} & \textbf{Requirement} \\ \midrule
\endhead
NFR-AUD-1 & The system shall maintain records of significant administrative and operational activities. \\
NFR-AUD-2 & Audit records shall identify the actor, action, affected record, and timestamp where applicable. \\
NFR-AUD-3 & Audit records shall be accessible only to authorized roles. \\
NFR-AUD-4 & Standard users shall not be able to edit or delete protected audit records. \\
NFR-AUD-5 & The system shall preserve status histories for major workflow entities. \\
NFR-AUD-6 & Audit information shall support investigation of failed, unauthorized, delayed, or disputed operations. \\
NFR-AUD-7 & Administrative reports should allow filtering of relevant activity records. \\
\bottomrule
\end{longtable}

\subsection{Portability Requirements}
\begin{longtable}{@{}p{2.6cm}p{10.8cm}@{}}
\caption{Portability Requirements} \label{tab:nfr-port} \\
\toprule \textbf{ID} & \textbf{Requirement} \\ \midrule
\endfirsthead
\toprule \textbf{ID} & \textbf{Requirement} \\ \midrule
\endhead
NFR-PORT-1 & The client application architecture shall support deployment on Android, Web, and Windows through Flutter. \\
NFR-PORT-2 & Shared application logic shall be reused across supported platforms where practical. \\
NFR-PORT-3 & Backend services shall remain independent of a single client platform. \\
NFR-PORT-4 & External services shall be accessed through documented APIs. \\
NFR-PORT-5 & The FastAPI backend shall be deployable to another compatible cloud provider if migration becomes necessary. \\
NFR-PORT-6 & Configuration values shall be externalized to simplify deployment across development, testing, and production environments. \\
NFR-PORT-7 & The AI model shall be stored and loaded in a format compatible with the selected deployment approach. \\
\bottomrule
\end{longtable}

\subsection{External-Service Resilience Requirements}
\begin{longtable}{@{}p{2.6cm}p{10.8cm}@{}}
\caption{External-Service Resilience Requirements} \label{tab:nfr-ext} \\
\toprule \textbf{ID} & \textbf{Requirement} \\ \midrule
\endfirsthead
\toprule \textbf{ID} & \textbf{Requirement} \\ \midrule
\endhead
NFR-EXT-1 & The system shall handle temporary failure of Firebase, Cloudinary, Google Maps Platform, Render, or other external services. \\
NFR-EXT-2 & The system shall display understandable error messages when an external integration is unavailable. \\
NFR-EXT-3 & Failed integration requests shall not be reported as successful. \\
NFR-EXT-4 & The system shall use suitable request timeouts where implemented. \\
NFR-EXT-5 & The system should support retry procedures for appropriate temporary failures. \\
NFR-EXT-6 & Sensitive external-service errors shall be logged for administrative review without being exposed fully to standard users. \\
NFR-EXT-7 & The system shall avoid permanent dependence on temporary files stored on the Render local filesystem. \\
NFR-EXT-8 & The system shall monitor external-service quotas and usage limits where applicable. \\
NFR-EXT-9 & Production deployment shall use suitable billing alerts, API restrictions, and service monitoring. \\
\bottomrule
\end{longtable}

\subsection{Legal and Ethical Requirements}
\begin{longtable}{@{}p{2.6cm}p{10.8cm}@{}}
\caption{Legal and Ethical Requirements} \label{tab:nfr-eth} \\
\toprule \textbf{ID} & \textbf{Requirement} \\ \midrule
\endfirsthead
\toprule \textbf{ID} & \textbf{Requirement} \\ \midrule
\endhead
NFR-ETH-1 & The system shall not present AI classification results as infallible. \\
NFR-ETH-2 & Users shall be allowed to correct or request review of inaccurate AI classifications. \\
NFR-ETH-3 & The system shall avoid discriminatory treatment based on user identity or location. \\
NFR-ETH-4 & Personal data shall be processed only for legitimate system functions. \\
NFR-ETH-5 & Electronic waste records shall not be used for unauthorized surveillance or unrelated purposes. \\
NFR-ETH-6 & Marketplace listings shall not knowingly support illegal, hazardous, counterfeit, or unauthorized goods. \\
NFR-ETH-7 & Administrative decisions affecting users shall be attributable to authorized personnel. \\
NFR-ETH-8 & Environmental and recycling reports shall not intentionally misrepresent incomplete or unverified information. \\
NFR-ETH-9 & The system shall distinguish implemented functionality from planned or future functionality in technical reporting. \\
\bottomrule
\end{longtable}

\subsection{Non-Functional Requirements Summary}

The non-functional requirements define the quality standards that EcoTrace
shall satisfy in addition to its operational functions. They address performance,
availability, reliability, security, privacy, usability, accessibility, scalability,
maintainability, compatibility, data integrity, backup, recovery, auditability,
external service resilience, and ethical operation.

These requirements shall guide the system architecture, database design,
interface design, implementation decisions, deployment configuration, and
testing activities. Chapter~\ref{ch:implementation} presents the test procedures
and available evidence used to evaluate whether the implemented system
satisfies these requirements -- including, per NFR-ETH-9 above, an explicit
account of which parts of the originally proposed technology stack were
delivered as designed and which were not.

\section{System Design}
\label{sec:system-design}

System design defines the technical structure, components, interfaces, data
models, and interactions required to implement the EcoTrace system. It
translates the functional and non-functional requirements into a structured
solution that can be developed, tested, deployed, and maintained.

EcoTrace was designed as a modular, cross-platform, cloud-based client-server
system. The design separates the user-interface layer, backend application layer,
data-management layer, external-service integration layer, and artificial
intelligence layer. This separation improves maintainability, scalability,
security, and testability because each component performs a clearly defined
responsibility and communicates with other components through controlled
interfaces.

The system design includes the overall system architecture, client
applications, backend services, cloud data management, authentication and
authorization, image-management services, location and route services,
notification services, artificial intelligence services, external interfaces, use
case models, data flow models, sequence diagrams, database design, and
user-interface design.

\subsection{System Architecture}
\label{sec:system-architecture}

The EcoTrace system follows a layered client-server architecture in which
Flutter client applications communicate with a FastAPI backend through
RESTful Application Programming Interfaces.

The architecture supports three primary client platforms: an Android mobile
application, a Web application, and a Windows desktop application. These
applications share a common Flutter codebase while providing interfaces
adapted to the operational requirements of different user roles and screen sizes.

The FastAPI backend acts as the central coordination layer. It processes client
requests, validates input, enforces business rules, verifies authentication and
authorization, communicates with external cloud services, performs AI
inference, and returns structured JSON responses.

Cloud Firestore stores structured system data, including user profiles,
electronic waste records, pickup requests, assignments, transportation records,
repair records, refurbishment records, recycling records, marketplace listings,
notifications, reports, and audit information.

Cloudinary stores and delivers uploaded media files, particularly electronic
waste images. Firebase Authentication manages user identity, while Firebase
Cloud Messaging delivers push notifications. Google Maps Platform provides
geocoding, location display, route calculation, nearby facility search, and route
optimization services.

The artificial intelligence component is implemented using PyTorch. It
analyses uploaded electronic waste images and returns a predicted waste
category and confidence score through the FastAPI backend.

The FastAPI backend is deployed on the Render cloud platform and is
accessed by the Flutter applications through secure HTTPS requests.

\subsubsection{Architectural Style}

EcoTrace combines several architectural principles.

\textbf{Client-Server Architecture.} The Flutter applications operate as
clients, while the FastAPI application provides centralized backend services.
Client applications submit requests and receive structured responses from the
server.

\textbf{Layered Architecture.} The system is divided into logical layers,
including presentation, application, integration, data, and infrastructure
layers. Each layer performs related responsibilities and reduces direct
dependency between unrelated components.

\textbf{RESTful Architecture.} Communication between Flutter and FastAPI
follows RESTful principles. Resources are accessed through HTTP endpoints,
while JSON is used as the primary data-exchange format.

\textbf{Cloud-Based Architecture.} Managed cloud services provide
authentication, database storage, notifications, media hosting, mapping, and
backend deployment.

\textbf{Modular Architecture.} System functions are divided into modules
such as authentication, waste registration, pickup management, repair,
recycling, notifications, and reporting. This makes the system easier to
maintain and extend.

\subsubsection{Presentation Layer}

The presentation layer consists of the Android, Web, and Windows
applications developed using Flutter. Its responsibilities include displaying
user interfaces, collecting user input, validating basic form data, displaying
loading, success, and error states, presenting dashboard information,
displaying maps, tables, charts, images, and notifications, enforcing
role-appropriate navigation, and communicating with the FastAPI backend and
approved Firebase services. The presentation layer does not make final
authorization decisions -- protected operations are verified again by the
backend.

\subsubsection{Application and Business Logic Layer}

The application and business logic layer is implemented using FastAPI. Its
responsibilities include processing API requests, validating request schemas,
verifying authentication tokens, enforcing role-based authorization, applying
workflow rules, generating identifiers and timestamps, coordinating pickup
assignments, processing image uploads, invoking the AI classification service,
communicating with Firebase Cloud Messaging, requesting Google Maps
services, creating reports and summaries, handling errors, and returning JSON
responses. Centralizing business logic in the backend reduces duplication and
prevents client applications from performing unauthorized operations.

\subsubsection{Data Management Layer}

The data-management layer uses Cloud Firestore as the primary database.
Cloud Firestore stores data as collections and documents; this document-oriented
model is suitable for EcoTrace because many records contain flexible
and nested attributes. The data-management layer is responsible for storing
structured records, retrieving records, updating workflow statuses, maintaining
relationships between documents, storing timestamps and ownership
information, supporting real-time synchronization, applying Firestore
Security Rules, and supporting administrative reporting and analytics. Media
files are not stored directly inside Firestore -- instead, Cloudinary stores the
files, while Firestore stores the corresponding secure URLs and metadata.

\subsubsection{Authentication and Authorization Layer}

Firebase Authentication provides user registration, login, password recovery,
session management, and identity tokens. After successful authentication, the
client receives an authentication token. Protected backend requests include the
token, which is verified before the requested operation is processed. Role and
profile information is stored in Cloud Firestore; the backend uses this
information to determine whether the authenticated user has permission to
perform the requested action. This layer supports the household, business,
collector, driver, technician, recycler, environmental officer, administrator, and
super administrator roles.

\subsubsection{Media Management Layer}

Cloudinary provides media storage, transformation, optimization, and
delivery. When a user uploads an electronic waste image: (1)~the Flutter
application sends the image to the FastAPI backend; (2)~the backend validates
the file; (3)~the backend uploads the file to Cloudinary; (4)~Cloudinary stores
and optimizes the image; (5)~Cloudinary returns a secure image URL; and
(6)~the backend stores the URL and related metadata in Cloud Firestore. This
approach prevents permanent media files from being stored on the ephemeral
Render filesystem.

\subsubsection{Notification Layer}

Firebase Cloud Messaging provides push-notification services. The
notification layer supports events such as account verification, pickup
submission, request approval, collector or driver assignment, pickup
reminders, collection completion, repair-status updates, refurbishment
completion, recycling completion, marketplace activity, and system
announcements. Notification records are stored in Cloud Firestore so users can
review their notification history within the application.

\subsubsection{Geospatial and Routing Layer}

Google Maps Platform provides the location-based services used by
EcoTrace. The layer supports displaying interactive maps, converting addresses
into coordinates, converting coordinates into addresses, recording pickup
locations, locating collection and processing centres, calculating distance and
estimated travel time, displaying suggested routes, and optimizing multiple
pickup locations where enabled. The FastAPI backend may process server-side
requests to Google Maps services and return only the required route or
location information to the client.

\subsubsection{Artificial Intelligence Layer}

The artificial intelligence layer is responsible for electronic waste image
classification. Its workflow includes: (1)~receiving an uploaded image;
(2)~validating the image; (3)~resizing and normalizing the image; (4)~loading
the trained PyTorch model; (5)~performing inference; (6)~identifying the
predicted class; (7)~calculating a confidence score; (8)~returning the result to
the backend; and (9)~storing the accepted result in Cloud Firestore. The model
architecture is designed to be replaceable so that a more accurate or efficient
model can be introduced without redesigning the complete system.

\subsubsection{Deployment and Infrastructure Layer}

The FastAPI backend is deployed as a Render Web Service. The deployment
layer provides a publicly accessible HTTPS endpoint, Git-based deployment,
environment-variable management, build and start commands, runtime logs,
deployment logs, and health-check support. The current free deployment is
suitable for testing and demonstration; however, the architecture permits
migration to a more powerful hosting plan or another compatible cloud
provider.

\subsubsection{System Architecture Data Flow}

The general data flow within the EcoTrace architecture is as follows:
(1)~a user performs an operation using an Android, Web, or Windows
application; (2)~the Flutter application validates the basic input; (3)~a secure
request is sent to the FastAPI backend; (4)~the backend verifies the user's
authentication token and role; (5)~the backend validates the submitted data;
(6)~the backend applies the relevant business rules; (7)~the backend
communicates with Cloud Firestore, Cloudinary, Firebase Cloud Messaging,
Google Maps Platform, or PyTorch where required; (8)~the external service
returns the requested result; (9)~the backend stores or updates the relevant
system records; (10)~the backend returns a structured JSON response; and
(11)~the Flutter application displays the result to the user.
Figure~\ref{fig:system-architecture} presents the major layers, components,
services, and communication paths within the EcoTrace system architecture, as
proposed at the design stage of this project.

\begin{figure}[H]
  \centering
  \begin{tikzpicture}[
      layer/.style={draw, rounded corners, minimum width=13.2cm, minimum height=1.15cm, align=center, font=\small},
      svc/.style={draw, rounded corners, fill=ecotracegreen!8, minimum width=2.6cm, minimum height=1cm, align=center, font=\scriptsize},
      node distance=6pt
    ]
    \node[layer, fill=ecotracegreen!18] (pres) {\textbf{Presentation Layer} -- Flutter (Android, Web, Windows)};
    \node[layer, below=of pres, fill=ecotracegreen!14] (app) {\textbf{Application \& Business Logic Layer} -- FastAPI (RESTful API)};
    \node[below=of app] (svcrow) {};
    \node[svc, below=10pt of app, xshift=-5.4cm] (db) {Cloud\\Firestore};
    \node[svc, right=6pt of db] (auth) {Firebase\\Authentication};
    \node[svc, right=6pt of auth] (media) {Cloudinary};
    \node[svc, right=6pt of media] (notif) {Firebase Cloud\\Messaging};
    \node[svc, right=6pt of notif] (maps) {Google Maps\\Platform};
    \node[svc, right=6pt of maps] (ai) {PyTorch\\AI Classifier};
    \node[layer, below=28pt of app, fill=ecotracegreen!10] (deploy) {\textbf{Deployment Layer} -- Render Cloud Platform (HTTPS, Render Web Service)};

    \draw[-Latex] (pres) -- (app);
    \draw[-Latex] (app.south) -- (svcrow.center);
    \foreach \s in {db,auth,media,notif,maps,ai} { \draw[-Latex] (app.south) -- (\s.north); }
    \draw[-Latex] (db.south) |- ($(deploy.north)+(0,4pt)$);
    \draw[-Latex] (app) -- (deploy);
  \end{tikzpicture}
  \caption{Overall System Architecture of EcoTrace, as proposed at the design stage. Source: Author's design.}
  \label{fig:system-architecture}
\end{figure}

\begin{table}[H]
\centering
\caption{Major Components of the EcoTrace System Architecture}
\label{tab:architecture-components}
\begin{tabular}{@{}p{2.6cm}p{3cm}p{7cm}@{}}
\toprule
\textbf{Architectural Layer} & \textbf{Technology} & \textbf{Primary Responsibility} \\
\midrule
Presentation Layer & Flutter (Android, Web, Windows) & Provides cross-platform user interfaces, navigation, forms, dashboards, maps, notifications, and role-specific interactions. \\
Application Layer & FastAPI & Processes requests, validates data, applies business rules, verifies authorization, coordinates integrations, and returns JSON responses. \\
Deployment Layer & Render & Hosts the FastAPI backend and provides HTTPS, environment variables, logging, health checks, and Git-based deployment. \\
Database Layer & Cloud Firestore & Stores user profiles, electronic waste records, workflows, notifications, reports, settings, and audit information. \\
Authentication Layer & Firebase Authentication & Provides account registration, login, password recovery, session management, and identity tokens. \\
Media Layer & Cloudinary & Stores, optimizes, transforms, and delivers electronic waste images. \\
Notification Layer & Firebase Cloud Messaging & Delivers targeted and system-wide push notifications. \\
Geospatial Layer & Google Maps Platform & Provides mapping, geocoding, location search, distance calculation, routing, and route optimization. \\
Artificial Intelligence Layer & PyTorch & Provides image preprocessing, electronic waste classification, and confidence-score generation. \\
\bottomrule
\end{tabular}
\end{table}

\subsection{Use Case Diagram and Actor Interactions}

A use case diagram represents the interactions between external actors and the
functions provided by a software system. It defines the system boundary,
identifies the categories of users, and shows the principal services available to
each actor.

The EcoTrace use case model contains ten main actors: (1)~Guest User;
(2)~Household User; (3)~Business User; (4)~Collector; (5)~Driver;
(6)~Technician; (7)~Recycler; (8)~Environmental Officer; (9)~Administrator;
and (10)~Super Administrator. Each actor performs different functions
according to assigned responsibilities and access permissions. Firebase
Authentication identifies registered users, while role-based access control
determines the functions that each actor can access.

\subsubsection{System Boundary}

The EcoTrace system boundary contains the major functions delivered
through the Android, Web, and Windows applications and the FastAPI backend.
External actors interact with the system through interfaces provided by the
presentation layer. The system processes requests, applies business rules,
communicates with external cloud services, stores records, and returns the
appropriate results.

The principal use case groups within the system boundary include
authentication and account management, electronic waste registration,
AI-assisted waste classification, pickup scheduling, collector and driver
assignment, collection and transportation, repair and refurbishment, recycling
and material recovery, marketplace management, notifications, environmental
monitoring, reporting and analytics, user and facility management, security
and audit management, and system configuration.

\subsubsection{Actor Interactions}

\textbf{Guest User.} The Guest User is an unauthenticated actor who accesses
public information and account-entry functions. The Guest User can view
public electronic waste awareness information, view publicly available
collection-centre information, register a new account, log in to an existing
account, and request password recovery. The Guest User must register and
authenticate before accessing protected EcoTrace functions.

\textbf{Household User.} The Household User represents an individual or
family that generates electronic waste. The Household User can manage a
personal profile, register electronic waste, upload waste images, request AI
classification, submit pickup requests, select a pickup location, track pickup
and collection status, receive notifications, view repair, refurbishment, and
recycling outcomes, browse marketplace listings, and view personal electronic
waste records.

\textbf{Business User.} The Business User represents an organization,
institution, company, or public office generating electronic waste. The
Business User can manage an organizational profile, register individual or
multiple electronic waste items, submit bulk collection requests, upload
images and supporting information, track organizational pickup requests,
monitor repair, recycling, and disposal outcomes, receive notifications, access
organizational reports, and maintain electronic waste disposal records.

\textbf{Collector.} The Collector is responsible for receiving and completing
pickup assignments. The Collector can view assigned pickup requests, accept
or reject assignments, view pickup locations and waste details, update
collection status, confirm item collection, upload collection evidence, report
collection exceptions, transfer collected waste to the next responsible actor,
and view completed collection history.

\textbf{Driver.} The Driver transports electronic waste between pickup
locations, collection centres, repair facilities, refurbishment facilities, and
recycling centres. The Driver can view assigned transportation tasks, access
pickup and destination locations, view suggested routes, follow optimized
routes where enabled, update departure, transit, arrival, and delivery status,
report delays or transportation incidents, record delivery details, and confirm
delivery to the receiving facility.

\textbf{Technician.} The Technician inspects, diagnoses, repairs, and
refurbishes electronic devices. The Technician can view assigned devices,
inspect and diagnose devices, determine repairability, record required parts
and estimated costs, update repair status, record replaced components,
perform functional testing, record refurbishment activities, assign quality
grades, approve devices for reuse or resale where authorized, and recommend
recycling for non-repairable devices.

\textbf{Recycler.} The Recycler processes electronic waste that cannot be
repaired or refurbished. The Recycler can view assigned recycling records,
receive and confirm electronic waste items, record dismantling activities,
record recovered components and materials, identify hazardous materials,
record residual waste, specify disposal methods, update recycling status,
confirm recycling completion, and generate recycling and material-recovery
reports.

\textbf{Environmental Officer.} The Environmental Officer monitors
electronic waste activities, environmental performance, and
compliance-related information. The Environmental Officer can view electronic
waste locations, monitor collection coverage, review transportation, repair,
refurbishment, and recycling activities, identify delayed or unresolved cases,
review material-recovery statistics, inspect facility records, record
environmental observations where implemented, view maps and analytical
dashboards, and generate environmental reports.

\textbf{Administrator.} The Administrator manages daily system operations
and coordinates the major EcoTrace workflows. The Administrator can manage
standard user accounts, verify operational users and organizations, manage
electronic waste categories, manage collection and processing facilities,
review pickup requests, assign collectors and drivers, monitor active
workflows, manage marketplace listings, publish announcements, review
operational reports, manage selected system content, and review permitted
audit records.

\textbf{Super Administrator.} The Super Administrator has the highest level
of system authority and manages high-risk administrative and security
operations. The Super Administrator can create and manage administrator
accounts, manage roles and permissions, configure system-wide settings,
review security and audit records, suspend or restore critical accounts, manage
integration configuration, review system health information, control
high-risk administrative functions, and monitor administrator activities.

\subsubsection{Major Use Case Relationships}

Some EcoTrace use cases depend on other use cases. These dependencies can
be represented using the \emph{include} relationship. Examples include:
Register Electronic Waste includes Validate Waste Information; Upload Waste
Image includes Validate Image File; Classify Electronic Waste includes
Preprocess Image; Submit Pickup Request includes Register Pickup Location;
Assign Collector includes Review Pickup Request; Complete Collection
includes Update Waste Lifecycle Status; Complete Repair includes Record
Repair Outcome; Complete Recycling includes Record Recovered Materials;
Generate Report includes Retrieve Authorized Data; and Perform Protected
Operation includes Verify Authentication and Role.

Optional or conditional behaviour may be represented using the
\emph{extend} relationship. Examples include: AI Classification extends
Register Electronic Waste; Upload Collection Evidence extends Complete
Collection; Report Transportation Incident extends Transport Electronic
Waste; Refurbish Device extends Repair Device; List Refurbished Item extends
Complete Refurbishment; and Send Status Notification extends major
workflow-status changes.

\subsubsection{Use Case Interaction Summary}

The use case model demonstrates that EcoTrace is a multi-actor system in
which each user category performs specialized functions. Household and
Business Users initiate the electronic waste lifecycle by registering items and
requesting collection. Collectors and Drivers coordinate collection and
transportation. Technicians determine whether devices can be repaired or
refurbished, while Recyclers process non-repairable waste and recover
materials. Environmental Officers monitor operational and environmental
performance. Administrators coordinate daily system activities, while Super
Administrators manage high-level security, permissions, and configuration.

The separation of actor responsibilities supports accountability, security,
workflow control, and complete electronic waste lifecycle traceability.
Figure~\ref{fig:use-case} presents the principal actors and use cases supported
by the EcoTrace system, and Table~\ref{tab:actor-summary} summarizes each
actor's principal use cases.

\begin{figure}[H]
  \centering
  \begin{tikzpicture}[
      actor/.style={draw, align=center, font=\scriptsize, minimum width=1.9cm},
      uc/.style={draw, ellipse, align=center, font=\scriptsize, fill=ecotracegreen!10, minimum width=2.6cm, minimum height=0.85cm},
      boundary/.style={draw, rounded corners, minimum width=8.6cm, minimum height=8.6cm}
    ]
    \node[boundary, label={[font=\small\bfseries]above:EcoTrace System Boundary}] (b) at (4.3,-4.3) {};
    \node[actor] (guest) at (-1.6,-0.4) {Guest\\User};
    \node[actor] (hh) at (-1.6,-1.6) {Household /\\Business User};
    \node[actor] (coll) at (-1.6,-2.8) {Collector /\\Driver};
    \node[actor] (tech) at (-1.6,-4.0) {Technician /\\Recycler};
    \node[actor] (env) at (-1.6,-5.2) {Environmental\\Officer};
    \node[actor] (admin) at (-1.6,-6.4) {Administrator /\\Super Admin};

    \node[uc] (auth) at (2.4,-0.7) {Register /\\Authenticate};
    \node[uc] (reg) at (5.6,-0.7) {Register \&\\Classify Waste};
    \node[uc] (pick) at (2.4,-2.0) {Pickup\\Request};
    \node[uc] (asgn) at (5.6,-2.0) {Assign \&\\Collect};
    \node[uc] (repr) at (2.4,-3.3) {Repair \&\\Refurbish};
    \node[uc] (recy) at (5.6,-3.3) {Recycle \&\\Recover};
    \node[uc] (mkt) at (2.4,-4.6) {Marketplace};
    \node[uc] (not) at (5.6,-4.6) {Notifications};
    \node[uc] (env2) at (2.4,-5.9) {Environmental\\Monitoring};
    \node[uc] (rpt) at (5.6,-5.9) {Reporting \&\\Analytics};
    \node[uc] (adm2) at (2.4,-7.2) {User \&\\Facility Mgmt};
    \node[uc] (sec) at (5.6,-7.2) {Security \&\\Audit};

    \draw (guest) -- (auth);
    \draw (hh) -- (auth); \draw (hh) -- (reg); \draw (hh) -- (pick); \draw (hh) -- (mkt); \draw (hh) -- (not);
    \draw (coll) -- (asgn); \draw (coll) -- (not);
    \draw (tech) -- (repr); \draw (tech) -- (recy); \draw (tech) -- (mkt);
    \draw (env) -- (env2); \draw (env) -- (rpt);
    \draw (admin) -- (adm2); \draw (admin) -- (sec); \draw (admin) -- (rpt); \draw (admin) -- (asgn);
  \end{tikzpicture}
  \caption{Overall Use Case Diagram for EcoTrace. Source: Author's design.}
  \label{fig:use-case}
\end{figure}

\begin{table}[H]
\centering
\small
\caption{Summary of EcoTrace Actors and Principal Use Cases}
\label{tab:actor-summary}
\begin{tabular}{@{}p{3.2cm}p{9.4cm}@{}}
\toprule
\textbf{Actor} & \textbf{Principal Use Cases} \\
\midrule
Guest User & View public information, register an account, log in, and request password recovery. \\
Household User & Register electronic waste, upload images, request AI classification, schedule pickups, track status, receive notifications, and browse the marketplace. \\
Business User & Manage organizational waste records, submit bulk pickup requests, monitor recovery outcomes, and generate organizational reports. \\
Collector & View assignments, accept pickup tasks, collect waste, upload evidence, update statuses, and complete collections. \\
Driver & View transport assignments, access routes, update transport status, report incidents, and confirm delivery. \\
Technician & Inspect devices, diagnose faults, repair devices, refurbish devices, record test results, and recommend recycling. \\
Recycler & Receive waste, dismantle devices, record recovered materials, identify hazardous waste, complete recycling, and generate recovery records. \\
Environmental Officer & Monitor activities, inspect maps and reports, review recovery performance, identify unresolved cases, and generate environmental reports. \\
Administrator & Manage users, categories, facilities, requests, assignments, marketplace listings, announcements, and operational reports. \\
Super Administrator & Manage administrators, roles, permissions, security, audit logs, integrations, and system-wide configuration. \\
\bottomrule
\end{tabular}
\end{table}

\textit{Note: the actor-to-use-case mapping in Table~\ref{tab:actor-summary} was
reconstructed from the source document's own actor descriptions, since the
original table's layout did not survive text extraction from the source PDF
cleanly. The use-case content itself is unchanged.}

\subsection{Data Flow Diagram}

A Data Flow Diagram (DFD) illustrates how information moves between
external entities, system processes, data stores, and outputs. It describes the
logical movement and transformation of data without specifying the
programming code or physical implementation used to perform each
operation. The EcoTrace data flow model is presented at two levels: the
Context-Level DFD, which represents EcoTrace as one complete process and
shows its interactions with external entities; and the Level-One DFD, which
decomposes the complete system into its principal operational processes and
data stores.

\subsubsection{Context-Level Data Flow Diagram}

The context-level DFD represents the complete EcoTrace system as a single
process identified as Process 0. It defines the system boundary and shows the
principal information exchanged between EcoTrace and its external users and
services. To maintain readability, actors with closely related responsibilities
are grouped within the context diagram; their detailed responsibilities are
expanded in the Level-One DFD and the use case model.

The principal external entities are: (1)~Guest and Public User;
(2)~Household and Business User; (3)~Collector and Driver; (4)~Technician
and Recycler; (5)~Environmental Officer; (6)~Administrator and Super
Administrator; and (7)~External Cloud and Artificial Intelligence Services
(Firebase Authentication, Cloud Firestore, Cloudinary, Firebase Cloud
Messaging, Google Maps Platform, PyTorch, and Render).
Figure~\ref{fig:context-dfd} presents the principal information flows between
EcoTrace, its users, and the external cloud services supporting the system, and
Table~\ref{tab:context-dfd} summarizes the data exchanged with each entity.

\begin{figure}[H]
  \centering
  \begin{tikzpicture}[
      ent/.style={draw, align=center, font=\scriptsize, minimum width=2.6cm, minimum height=0.9cm},
      proc/.style={draw, circle, align=center, font=\small\bfseries, fill=ecotracegreen!14, minimum size=2.6cm},
      arr/.style={-Latex, font=\tiny}
    ]
    \node[proc] (p0) at (0,0) {Process 0\\EcoTrace};
    \node[ent] (a) at (-5.2,2.4) {Guest \&\\Public User};
    \node[ent] (b) at (-5.2,0.8) {Household \&\\Business User};
    \node[ent] (c) at (-5.2,-0.8) {Collector \&\\Driver};
    \node[ent] (d) at (-5.2,-2.4) {Technician \&\\Recycler};
    \node[ent] (e) at (5.2,2.4) {Environmental\\Officer};
    \node[ent] (f) at (5.2,0.8) {Administrator \&\\Super Admin};
    \node[ent] (g) at (5.2,-1.6) {External Cloud\\\& AI Services};
    \draw[arr] (a) -- (p0); \draw[arr] (b) -- (p0); \draw[arr] (c) -- (p0); \draw[arr] (d) -- (p0);
    \draw[arr] (p0) -- (e); \draw[arr] (p0) -- (f); \draw[arr] (p0) -- (g);
  \end{tikzpicture}
  \caption{Context-Level Data Flow Diagram of the EcoTrace System. Source: Author's design.}
  \label{fig:context-dfd}
\end{figure}

\begin{longtable}{@{}p{3.1cm}p{5cm}p{5.3cm}@{}}
\caption{External Entities and Data Flows in the EcoTrace Context Diagram} \label{tab:context-dfd} \\
\toprule \textbf{External Entity} & \textbf{Data Sent to EcoTrace} & \textbf{Data Received from EcoTrace} \\ \midrule
\endfirsthead
\toprule \textbf{External Entity} & \textbf{Data Sent to EcoTrace} & \textbf{Data Received from EcoTrace} \\ \midrule
\endhead
Guest and Public User & Registration data, login credentials, password-recovery requests, public-information requests. & Public information, registration confirmation, authentication results, recovery instructions. \\
Household and Business User & Profile data, electronic waste records, images, pickup requests, locations, marketplace enquiries, report requests. & Classification results, pickup status, collection updates, recovery outcomes, notifications, marketplace information, reports. \\
Collector and Driver & Task responses, collection updates, transport updates, evidence, delivery confirmations, incident reports. & Assignments, waste details, schedules, pickup locations, destinations, route information. \\
Technician and Recycler & Inspection findings, repair records, refurbishment records, recycling records, material-recovery information. & Assigned items, device details, lifecycle history, processing instructions. \\
Environmental Officer & Monitoring requests, report filters, inspection observations, compliance findings. & Maps, environmental reports, alerts, recovery statistics, facility records, analytical dashboards. \\
Administrator and Super Administrator & User-management actions, assignments, role changes, facility data, announcements, report requests, configuration changes. & Administrative dashboards, user and account records, workflow status, assignment results, system reports, audit logs, security alerts, configuration status. \\
External Cloud and AI Services & Authentication requests, database operations, media files, notification requests, geographic requests, images for AI classification. & Authentication tokens, stored data, image URLs, delivery results, routes, location information, prediction results. \\
\bottomrule
\end{longtable}

\subsubsection{Level-One Data Flow Diagram}

The Level-One DFD expands Process 0 into its major internal processes,
data stores, and information flows. It provides a more detailed logical view of
how data moves through the system during registration, collection, logistics,
recovery, reporting, and administration. The system is decomposed into eight
major processes, each reading from and writing to its own principal data
store: 1.0~Manage Authentication and User Accounts (D1); 2.0~Manage
Electronic Waste Registration and AI Classification (D2); 3.0~Manage Pickup
Requests and Assignments (D3); 4.0~Manage Collection, Transportation, and
Routing (D4); 5.0~Manage Repair, Refurbishment, and Recycling (D5);
6.0~Manage Marketplace and Notifications (D6); 7.0~Manage Reporting and
Environmental Monitoring (D7); and 8.0~Manage Administration, Security,
and Configuration (D8).

\begin{itemize}
  \item \textbf{Process 1.0} handles registration, login, password recovery, and
    profile management. It communicates with Firebase Authentication for
    identity verification and reads/writes D1~--~User Accounts and Profiles.
  \item \textbf{Process 2.0} manages waste registration, image upload, and
    AI-based classification. It writes to D2~--~Electronic Waste Records,
    uploads images through Cloudinary, and communicates with the PyTorch
    classification service.
  \item \textbf{Process 3.0} handles pickup scheduling, validation, review, and
    collector/driver assignment. It reads D2 and writes D3~--~Pickup and
    Assignment Records.
  \item \textbf{Process 4.0} manages collection execution, transport, route
    display, and delivery confirmation. It communicates with Google Maps
    Platform and writes D4~--~Transportation and Route Records, updating D3.
  \item \textbf{Process 5.0} manages inspection, repair, refurbishment,
    recycling, and material recovery. It reads D2 and writes D5~--~Repair,
    Refurbishment, and Recycling Records.
  \item \textbf{Process 6.0} manages marketplace listings and notification
    delivery. It writes D6~--~Marketplace and Notification Records and
    communicates with Firebase Cloud Messaging.
  \item \textbf{Process 7.0} generates dashboards, reports, and monitoring
    views. It reads D2--D6 and writes D7~--~Reports, Monitoring, and Audit
    Records.
  \item \textbf{Process 8.0} manages user administration, roles, facilities,
    audit review, and configuration. It updates D1, D7, and D8~--~Configuration
    and Reference Data.
\end{itemize}

This decomposition shows that authentication, waste registration, pickup
management, logistics, recovery operations, reporting, monitoring, and
administration are handled through distinct but interconnected processes, and
that Cloud Firestore-based records, media services, notification services,
geospatial services, and AI services are integrated into the logical operation of
the system. Figure~\ref{fig:level1-dfd} illustrates the major internal processes,
data stores, and data flows, and Table~\ref{tab:level1-dfd} summarizes each
process.

\begin{figure}[H]
  \centering
  \begin{tikzpicture}[
      proc/.style={draw, rounded corners, align=center, font=\scriptsize, fill=ecotracegreen!12, minimum width=2.5cm, minimum height=0.85cm},
      store/.style={draw, minimum width=2.3cm, minimum height=0.6cm, font=\scriptsize, align=center},
      arr/.style={-Latex, font=\tiny}
    ]
    \node[proc] (p1) at (0,3.2)   {1.0 Auth \&\\User Accounts};
    \node[proc] (p2) at (3.4,3.2) {2.0 Waste Reg.\\\& AI Classify};
    \node[proc] (p3) at (6.8,3.2) {3.0 Pickup \&\\Assignment};
    \node[proc] (p4) at (0,0)     {4.0 Collection,\\Transport, Routing};
    \node[proc] (p5) at (3.4,0)   {5.0 Repair, Refurb.,\\Recycling};
    \node[proc] (p6) at (6.8,0)   {6.0 Marketplace \&\\Notifications};
    \node[proc] (p7) at (1.7,-3.2) {7.0 Reporting \&\\Env. Monitoring};
    \node[proc] (p8) at (5.1,-3.2) {8.0 Admin, Security,\\Configuration};
    \node[store] (d1) at (-2.6,3.2) {D1 Users};
    \node[store] (d2) at (3.4,1.9) {D2 E-waste};
    \node[store] (d3) at (9.6,3.2) {D3 Pickups};
    \node[store] (d4) at (-2.6,0) {D4 Transport};
    \node[store] (d5) at (3.4,-1.3) {D5 Repair/Recycle};
    \node[store] (d6) at (9.6,0) {D6 Marketplace};
    \node[store] (d7) at (1.7,-4.6) {D7 Reports/Audit};
    \node[store] (d8) at (5.1,-4.6) {D8 Config};
    \foreach \p/\d in {p1/d1,p2/d2,p3/d3,p4/d4,p5/d5,p6/d6,p7/d7,p8/d8} { \draw[arr] (\p) -- (\d); }
  \end{tikzpicture}
  \caption{Level-One Data Flow Diagram of the EcoTrace System. Source: Author's design.}
  \label{fig:level1-dfd}
\end{figure}

\begin{longtable}{@{}p{1.4cm}p{4.4cm}p{6.6cm}@{}}
\caption{Major Processes in the EcoTrace Level-One Data Flow Diagram} \label{tab:level1-dfd} \\
\toprule \textbf{Process} & \textbf{Process Name} & \textbf{Primary Function} \\ \midrule
\endfirsthead
\toprule \textbf{Process} & \textbf{Process Name} & \textbf{Primary Function} \\ \midrule
\endhead
1.0 & Manage Authentication and User Accounts & Handles registration, login, password recovery, profile management, role retrieval, and access verification. \\
2.0 & Manage Electronic Waste Registration and AI Classification & Handles waste registration, image upload, item storage, and AI-based classification of electronic waste images. \\
3.0 & Manage Pickup Requests and Assignments & Handles pickup scheduling, request validation, review, and assignment of collectors and drivers. \\
4.0 & Manage Collection, Transportation, and Routing & Handles collection operations, transportation tracking, route display, delivery updates, and related logistics. \\
5.0 & Manage Repair, Refurbishment, and Recycling & Handles inspection, diagnosis, repair, refurbishment, recycling, and material-recovery activities. \\
6.0 & Manage Marketplace and Notifications & Handles marketplace listings, browsing, announcements, and push notifications. \\
7.0 & Manage Reporting and Environmental Monitoring & Handles dashboards, analytical views, reports, environmental monitoring, and related queries. \\
8.0 & Manage Administration, Security, and Configuration & Handles user administration, roles, facilities, audit review, security oversight, and system settings. \\
\bottomrule
\end{longtable}

\subsection{Sequence Diagrams}

A sequence diagram shows how actors and system components interact over
time to complete a specific operation. It emphasizes the order of messages
exchanged between the user interface, backend services, databases, and
external platforms. The principal EcoTrace sequence diagrams cover user
authentication, electronic waste registration and image upload, AI-based
classification, pickup request and assignment, collection and transportation,
repair/refurbishment/recycling, and report generation and environmental
monitoring.

\subsubsection{User Authentication Sequence}
The sequence proceeds as follows: (1)~the user enters an email address and
password; (2)~the Flutter application validates the required fields; (3)~the
application submits the credentials to Firebase Authentication; (4)~Firebase
Authentication verifies the credentials; (5)~if valid, Firebase returns an
authentication token; (6)~the Flutter application sends the token to the
FastAPI backend; (7)~FastAPI verifies the token and retrieves the user's
profile and role from Cloud Firestore; (8)~Cloud Firestore returns the profile
and role information; (9)~FastAPI returns an authorized session response; and
(10)~the Flutter application redirects the user to the appropriate role-based
dashboard. If authentication fails, Firebase Authentication returns an error
and the Flutter application displays an appropriate message.

\begin{figure}[H]
  \centering
  \begin{tikzpicture}[font=\scriptsize]
    \foreach \name/\x in {User/0,Flutter/2.4,Firebase/4.8,FastAPI/7.2,Firestore/9.6} {
      \node at (\x,0) {\textbf{\name}};
      \draw[dashed] (\x,-0.3) -- (\x,-6.3);
    }
    \draw[-Latex] (0,-1) -- (2.4,-1) node[midway,above]{credentials};
    \draw[-Latex] (2.4,-1.7) -- (4.8,-1.7) node[midway,above]{submit login};
    \draw[-Latex] (4.8,-2.4) -- (2.4,-2.4) node[midway,above]{auth token};
    \draw[-Latex] (2.4,-3.1) -- (7.2,-3.1) node[midway,above]{token};
    \draw[-Latex] (7.2,-3.8) -- (9.6,-3.8) node[midway,above]{get profile/role};
    \draw[-Latex] (9.6,-4.5) -- (7.2,-4.5) node[midway,above]{profile/role};
    \draw[-Latex] (7.2,-5.2) -- (2.4,-5.2) node[midway,above]{session response};
    \draw[-Latex] (2.4,-5.9) -- (0,-5.9) node[midway,above]{role dashboard};
  \end{tikzpicture}
  \caption{Authentication Sequence Diagram for EcoTrace. Source: Author's design.}
  \label{fig:seq-auth}
\end{figure}

\subsubsection{Electronic Waste Registration and Image Upload Sequence}
(1)~The user opens the electronic waste registration form; (2)~the Flutter
application displays the required fields; (3)~the user enters waste details and
selects an image; (4)~the Flutter application validates the basic input; (5)~the
application sends the image and registration data to the FastAPI backend;
(6)~FastAPI verifies the user's authentication token and role; (7)~FastAPI
validates the submitted data and image file; (8)~FastAPI uploads the image to
Cloudinary; (9)~Cloudinary stores the image and returns a secure image URL;
(10)~FastAPI creates the electronic waste record in Cloud Firestore, including
the image URL and lifecycle status; (11)~Cloud Firestore returns the created
record identifier; (12)~FastAPI returns a successful registration response; and
(13)~the Flutter application displays the registered item and confirmation
message.

\subsubsection{AI Classification Sequence}
(1)~The user selects the AI classification function; (2)~the Flutter
application sends the image or image reference to the FastAPI backend;
(3)~FastAPI verifies the user's authentication and authorization; (4)~FastAPI
retrieves or receives the image; (5)~the backend validates and preprocesses the
image; (6)~the preprocessed image is passed to the trained PyTorch model;
(7)~the PyTorch model performs inference; (8)~the model returns the
predicted waste category and confidence score; (9)~FastAPI formats the
result; (10)~the result is stored with the relevant waste record in Cloud
Firestore where appropriate; (11)~FastAPI returns the classification result to
the Flutter application; (12)~the Flutter application displays the predicted
category and confidence score; and (13)~the user may accept or correct the
result.

\subsubsection{Pickup Request and Assignment Sequence}
(1)~The user selects one or more registered electronic waste items; (2)~the
user enters a preferred pickup date, location, and collection instructions;
(3)~the Flutter application sends the pickup request to FastAPI; (4)~FastAPI
validates the request and user ownership; (5)~FastAPI stores the request in
Cloud Firestore; (6)~Cloud Firestore returns the request identifier; (7)~FastAPI
sends a notification request through Firebase Cloud Messaging; (8)~the
administrator receives a new-request notification; (9)~the administrator opens
the request and reviews its details; (10)~the administrator selects a collector
and driver; (11)~the Flutter administrative interface sends the assignment to
FastAPI; (12)~FastAPI validates the selected operational users; (13)~FastAPI
updates the pickup and assignment records in Cloud Firestore; (14)~FastAPI
sends assignment notifications through Firebase Cloud Messaging; and
(15)~the household or business user, collector, and driver receive updated
status information.
\begin{figure}[H]
  \centering
  \begin{tikzpicture}[font=\scriptsize]
    \foreach \name/\x in {User/0,Flutter/2.3,FastAPI/4.6,Firestore/6.9,FCM/9.2,Admin/11.3} {
      \node at (\x,0) {\textbf{\name}};
      \draw[dashed] (\x,-0.3) -- (\x,-5.3);
    }
    \draw[-Latex] (0,-0.9) -- (2.3,-0.9) node[midway,above,font=\tiny]{select items, date, location};
    \draw[-Latex] (2.3,-1.5) -- (4.6,-1.5) node[midway,above,font=\tiny]{submit request};
    \draw[-Latex] (4.6,-2.1) -- (6.9,-2.1) node[midway,above,font=\tiny]{store request};
    \draw[-Latex] (4.6,-2.7) -- (9.2,-2.7) node[midway,above,font=\tiny]{notify};
    \draw[-Latex] (9.2,-3.3) -- (11.3,-3.3) node[midway,above,font=\tiny]{new request};
    \draw[-Latex] (11.3,-3.9) -- (4.6,-3.9) node[midway,above,font=\tiny]{assign collector/driver};
    \draw[-Latex] (4.6,-4.5) -- (6.9,-4.5) node[midway,above,font=\tiny]{update assignment};
    \draw[-Latex] (4.6,-5.0) -- (9.2,-5.0) node[midway,above,font=\tiny]{notify all parties};
  \end{tikzpicture}
  \caption{Pickup, Assignment, Collection, and Transportation Sequence Diagram. Source: Author's design.}
  \label{fig:seq-pickup}
\end{figure}

\subsubsection{Collection and Transportation Sequence}
(1)~The collector opens the assigned pickup task; (2)~the Flutter application
requests the pickup details from FastAPI; (3)~FastAPI retrieves the request
and location information from Cloud Firestore; (4)~FastAPI requests route
information from Google Maps Platform where required; (5)~Google Maps
Platform returns distance, route, and estimated travel information; (6)~FastAPI
returns the task and route details to the collector or driver; (7)~the collector
accepts the task and updates the status; (8)~FastAPI stores the updated status
in Cloud Firestore; (9)~the collector reaches the pickup location and confirms
collection; (10)~the collector uploads evidence where required; (11)~FastAPI
stores the collection record and updates the waste lifecycle status; (12)~the
driver transports the waste to the assigned destination; (13)~the driver
updates transit and delivery status; (14)~FastAPI stores the transportation
record; (15)~the receiving facility confirms delivery; and (16)~Firebase Cloud
Messaging sends status notifications to the relevant users.

\subsubsection{Repair, Refurbishment, and Recycling Sequence}
(1)~The technician or recycler opens the assigned item; (2)~the Flutter
application requests the item and lifecycle history; (3)~FastAPI retrieves the
item, collection, and transport records from Cloud Firestore; (4)~the
technician inspects the device and records findings; (5)~FastAPI stores the
inspection record; (6)~if the item is repairable, the technician records repair
or refurbishment activities; (7)~FastAPI updates the relevant records in Cloud
Firestore; (8)~if the item is not repairable, it is assigned to the recycler;
(9)~the recycler records dismantling, recovered materials, hazardous
materials, and residual waste; (10)~FastAPI stores the recycling record and
updates the waste lifecycle status; (11)~where applicable, an approved
refurbished item is added to the marketplace; (12)~Firebase Cloud Messaging
sends completion or status notifications; and (13)~relevant users can view the
final recovery outcome.
\begin{figure}[H]
  \centering
  \begin{tikzpicture}[font=\scriptsize]
    \foreach \name/\x in {Technician/0,Flutter/2.6,FastAPI/5.2,Firestore/7.8,FCM/10.4} {
      \node at (\x,0) {\textbf{\name}};
      \draw[dashed] (\x,-0.3) -- (\x,-5.3);
    }
    \draw[-Latex] (0,-0.9) -- (2.6,-0.9) node[midway,above,font=\tiny]{open item};
    \draw[-Latex] (2.6,-1.5) -- (5.2,-1.5) node[midway,above,font=\tiny]{request history};
    \draw[-Latex] (5.2,-2.1) -- (7.8,-2.1) node[midway,above,font=\tiny]{fetch records};
    \draw[-Latex] (0,-2.7) -- (2.6,-2.7) node[midway,above,font=\tiny]{inspect / repair / recycle};
    \draw[-Latex] (2.6,-3.3) -- (5.2,-3.3) node[midway,above,font=\tiny]{submit outcome};
    \draw[-Latex] (5.2,-3.9) -- (7.8,-3.9) node[midway,above,font=\tiny]{update record};
    \draw[-Latex] (5.2,-4.5) -- (10.4,-4.5) node[midway,above,font=\tiny]{notify};
  \end{tikzpicture}
  \caption{Repair, Refurbishment, Recycling, and Reporting Sequence Diagram. Source: Author's design.}
  \label{fig:seq-repair}
\end{figure}

\subsubsection{Report Generation and Environmental Monitoring Sequence}
(1)~The authorized user opens the reporting dashboard; (2)~the user selects
report type, date range, status, category, location, or other filters; (3)~the
Flutter application sends the report request to FastAPI; (4)~FastAPI verifies
the user's role and report permissions; (5)~FastAPI retrieves relevant data
from Cloud Firestore; (6)~FastAPI aggregates and analyses the retrieved
records; (7)~FastAPI may request map or location information from Google
Maps Platform; (8)~FastAPI prepares the report data; (9)~the report is
returned to the Flutter application; (10)~the application displays tables,
charts, maps, and summary indicators; (11)~the authorized user may export
the report in PDF or another supported format; and (12)~the report request or
export action may be recorded in the audit log.

\subsubsection{Sequence Diagram Summary}
The sequence diagrams demonstrate how EcoTrace coordinates user actions,
Flutter applications, FastAPI services, Cloud Firestore, Firebase, Cloudinary,
Google Maps Platform, and PyTorch. They also show that authentication,
validation, authorization, cloud storage, notification delivery, route services,
and AI inference occur through ordered interactions rather than isolated
system operations.

\subsection{Database Design and Cloud Firestore Collections}

Database design defines how application data is structured, stored, related,
retrieved, updated, secured, and maintained. EcoTrace uses Cloud Firestore as
its primary database because it supports flexible document-based data
storage, real-time synchronization, cloud accessibility, security rules, and
integration with Firebase Authentication.

Cloud Firestore is a NoSQL document database. Instead of storing data in
relational tables, it organizes information into collections and documents. A
collection contains related documents, while each document contains fields
represented as key-value pairs. The EcoTrace database design separates
information according to the major system modules. This modular structure
improves maintainability, security, query performance, and future scalability.

The principal collections proposed for EcoTrace are: users, roles,
ewaste\_records, pickup\_requests, assignments, transportation\_records,
collection\_centres, facilities, repair\_records, refurbishment\_records,
recycling\_records, recovered\_materials, marketplace\_listings, notifications,
environmental\_reports, audit\_logs, system\_settings, and reference\_data.

\subsubsection{Collection Descriptions}

\textbf{users.} Stores profile and authorization-related information for
registered users. Firebase Authentication manages the user's email, password
credentials, authentication token, and authentication identifier, while Cloud
Firestore stores additional application-specific profile information. A typical
document may contain: userId, firebaseUid, fullName, email, phoneNumber,
roleId, organizationName, profileImageUrl, address, district, latitude,
longitude, accountStatus, isVerified, createdAt, updatedAt, and lastLoginAt.
The document identifier may use the Firebase Authentication user identifier
to simplify user-profile retrieval.

\textbf{roles.} Stores the available system roles and their associated
permissions: roleId, roleName, description, permissions, isActive, createdAt,
updatedAt. Examples of role values include household, business, collector,
driver, technician, recycler, environmental officer, administrator, and super
administrator.

\textbf{ewaste\_records.} Stores the main information associated with
registered electronic waste items: ewasteId, ownerId, ownerType, categoryId,
deviceName, brand, model, serialNumber, quantity, condition, description,
imageUrls, aiPredictedCategory, aiConfidenceScore, classificationSource,
latitude, longitude, address, preferredRecoveryOption, lifecycleStatus,
createdAt, updatedAt. The lifecycleStatus field may contain values such as
registered, pickup requested, assigned, collected, transported, under
inspection, under repair, refurbished, recycled, listed, reused, or disposed.

\textbf{pickup\_requests.} Stores requests submitted by household and
business users: pickupRequestId, requesterId, ewasteIds, pickupAddress,
latitude, longitude, preferredDate, preferredTime, contactName, contactPhone,
collectionInstructions, status, rejectionReason, cancellationReason,
submittedAt, scheduledAt, completedAt, updatedAt.

\textbf{assignments.} Stores the assignment of collectors, drivers,
technicians, and recyclers to operational records: assignmentId,
referenceType, referenceId, assignedUserId, assignedRole, assignedBy,
assignmentStatus, acceptedAt, rejectedAt, rejectionReason, assignedAt,
updatedAt. The referenceType field identifies whether the assignment relates
to pickup, transportation, repair, refurbishment, or recycling.

\textbf{transportation\_records.} Stores the movement of electronic waste
between operational locations: transportId, pickupRequestId, driverId,
vehicleId, originAddress, originLatitude, originLongitude, destinationAddress,
destinationLatitude, destinationLongitude, routeDistance, estimatedDuration,
routePolyline, status, incidentDescription, departedAt, arrivedAt, deliveredAt,
updatedAt.

\textbf{collection\_centres.} Stores approved electronic waste collection
locations: centreId, centreName, centreType, address, district, latitude,
longitude, contactPhone, email, operatingHours, acceptedWasteCategories,
verificationStatus, isActive, createdAt, updatedAt.

\textbf{facilities.} Stores repair centres, refurbishment centres, recycling
facilities, warehouses, and other processing locations: facilityId,
facilityName, facilityType, operatorId, address, latitude, longitude,
contactInformation, supportedOperations, capacity, verificationStatus,
isActive, createdAt, updatedAt.

\textbf{repair\_records.} Stores inspection, diagnosis, repair planning, and
repair-outcome information: repairId, ewasteId, technicianId, facilityId,
inspectionFindings, diagnosis, isRepairable, requiredParts, estimatedCost,
actualCost, replacedComponents, repairStatus, testResults, finalOutcome,
startedAt, completedAt, updatedAt.

\textbf{refurbishment\_records.} Stores information about devices restored
for reuse or resale: refurbishmentId, ewasteId, technicianId, facilityId,
activitiesPerformed, componentsReplaced, softwareUpdated,
functionalTestResults, cosmeticCondition, qualityGrade, approvalStatus,
approvedBy, finalCondition, startedAt, completedAt, updatedAt.

\textbf{recycling\_records.} Stores dismantling, processing,
material-recovery, and final-disposal information: recyclingId, ewasteId,
recyclerId, facilityId, quantityReceived, receivedCondition,
dismantlingActivities, hazardousMaterials, residualWaste, disposalMethod,
recyclingStatus, receivedAt, completedAt, updatedAt.

\textbf{recovered\_materials.} Stores materials and components recovered
during recycling: materialRecordId, recyclingId, materialType, materialName,
quantity, unit, qualityCondition, estimatedValue, destination,
marketplaceEligible, recordedBy, recordedAt.

\textbf{marketplace\_listings.} Stores refurbished devices, reusable
components, and recovered materials offered through the marketplace:
listingId, sourceRecordType, sourceRecordId, sellerId, title, description,
category, condition, qualityGrade, price, currency, quantityAvailable,
imageUrls, listingStatus, approvalStatus, approvedBy, createdAt, updatedAt.

\textbf{notifications.} Stores in-application notification records:
notificationId, recipientId, recipientRole, title, message, notificationType,
referenceType, referenceId, isRead, deliveryStatus, createdAt, readAt.

\textbf{environmental\_reports.} Stores submitted observations, inspection
findings, environmental assessments, and generated report metadata:
reportId, reportType, generatedBy, facilityId, location, reportPeriodStart,
reportPeriodEnd, observations, findings, recommendations, severity, status,
attachmentUrls, createdAt, updatedAt.

\textbf{audit\_logs.} Stores evidence of significant system and
administrative actions: auditId, actorId, actorRole, action, referenceType,
referenceId, previousValue, newValue, ipAddress, deviceInformation, result,
description, createdAt. Access to this collection shall be restricted to
authorized administrators and super administrators.

\textbf{system\_settings.} Stores configurable application settings: settingId,
settingName, settingValue, category, description, isSensitive, updatedBy,
updatedAt. Sensitive secrets shall not be stored directly in Cloud Firestore
when they can be managed through secure environment variables.

\textbf{reference\_data.} Stores controlled values used throughout the
system, such as electronic waste categories, account statuses, pickup statuses,
lifecycle statuses, facility types, material types, quality grades, notification
types, report types, and districts or operational regions. Centralizing these
values improves consistency across Android, Web, Windows, backend, and
reporting components.

\subsubsection{Database Relationships}

Although Cloud Firestore is a NoSQL database, logical relationships exist
between the collections: one user may register multiple electronic waste
records; one electronic waste record may belong to one owner; one pickup
request may contain multiple electronic waste records; one pickup request
may have one or more assignment records; one driver may complete multiple
transportation records; one electronic waste record may have repair,
refurbishment, or recycling records; one recycling record may contain
multiple recovered-material records; one refurbishment or recovered-material
record may generate a marketplace listing; one user may receive multiple
notifications; and one actor may generate multiple audit-log entries. These
relationships are maintained using document identifiers such as userId,
ewasteId, pickupRequestId, assignmentId, repairId, and recyclingId.
Figure~\ref{fig:er-diagram} presents these relationships as an entity
relationship diagram.

\begin{figure}[H]
  \centering
  \begin{tikzpicture}[
      ent/.style={draw, rounded corners, align=center, font=\scriptsize, fill=ecotracegreen!10, minimum width=2.5cm, minimum height=0.7cm},
      rel/.style={font=\tiny, midway, fill=white, inner sep=1pt}
    ]
    \node[ent] (users) at (0,4) {users};
    \node[ent] (ewaste) at (0,2) {ewaste\_records};
    \node[ent] (pickup) at (0,0) {pickup\_requests};
    \node[ent] (asgn) at (-3.4,0) {assignments};
    \node[ent] (trans) at (-3.4,-2) {transportation\_records};
    \node[ent] (repair) at (3.4,0) {repair\_records};
    \node[ent] (refurb) at (3.4,-2) {refurbishment\_records};
    \node[ent] (recyc) at (0,-2) {recycling\_records};
    \node[ent] (mat) at (0,-4) {recovered\_materials};
    \node[ent] (mkt) at (3.4,-4) {marketplace\_listings};
    \node[ent] (notif) at (-3.4,2) {notifications};
    \node[ent] (audit) at (3.4,2) {audit\_logs};
    \draw (users) -- (ewaste) node[rel]{1:N};
    \draw (users) -- (notif) node[rel]{1:N};
    \draw (users) -- (audit) node[rel]{1:N};
    \draw (ewaste) -- (pickup) node[rel]{N:M};
    \draw (pickup) -- (asgn) node[rel]{1:N};
    \draw (asgn) -- (trans) node[rel]{1:N};
    \draw (ewaste) -- (repair) node[rel]{1:N};
    \draw (repair) -- (refurb) node[rel]{1:1};
    \draw (ewaste) -- (recyc) node[rel]{1:N};
    \draw (recyc) -- (mat) node[rel]{1:N};
    \draw (mat) -- (mkt) node[rel]{1:N};
    \draw (refurb) -- (mkt) node[rel]{1:N};
  \end{tikzpicture}
  \caption{Cloud Firestore Database Schema (Entity Relationship Diagram) for EcoTrace. Source: Author's design.}
  \label{fig:er-diagram}
\end{figure}

\subsubsection{Indexing and Query Considerations}

Cloud Firestore queries require suitable indexing to support efficient data
retrieval. Likely indexed fields include user role and account status, owner
identifier and waste lifecycle status, pickup-request status and preferred date,
assigned user and assignment status, driver identifier and transportation
status, technician identifier and repair status, recycler identifier and recycling
status, listing category and listing status, notification recipient and read
status, report type and creation date, and audit actor and action date.
Composite indexes may be required when filtering and ordering records by
multiple fields.

\subsubsection{Database Security}

Cloud Firestore Security Rules shall restrict access according to
authentication status, ownership, role, and operation type. The security design
shall ensure that users can access only permitted personal records; users
cannot modify another person's private records; collectors and drivers access
only assigned operational records; technicians and recyclers access only
authorized processing records; environmental officers access approved
monitoring and reporting information; administrators access operational data
according to assigned permissions; super administrators access high-level
system data; and audit and configuration records remain protected. Backend
authorization checks shall complement Firestore Security Rules.

\subsubsection{Database Design Summary}

The database design provides a structured foundation for storing and
linking all major EcoTrace records. The proposed collections support
authentication, electronic waste registration, pickup scheduling, logistics,
repair, refurbishment, recycling, marketplace operations, notifications,
reporting, environmental monitoring, audit tracking, and configuration
management. The document-oriented structure provides flexibility while the
use of consistent identifiers, timestamps, statuses, security rules, and indexes
supports traceability, security, and efficient data retrieval.
Table~\ref{tab:collections-summary} summarizes the principal Cloud Firestore
collections and their representative fields.

\begin{longtable}{@{}p{3.2cm}p{4.6cm}p{5.4cm}@{}}
\caption{Summary of Major Cloud Firestore Collections} \label{tab:collections-summary} \\
\toprule \textbf{Collection} & \textbf{Purpose} & \textbf{Representative Fields} \\ \midrule
\endfirsthead
\toprule \textbf{Collection} & \textbf{Purpose} & \textbf{Representative Fields} \\ \midrule
\endhead
users & Stores user profiles, roles, account status, contact information, and location data. & userId, fullName, email, roleId, accountStatus, createdAt \\
ewaste\_records & Stores registered electronic waste details, images, AI results, and lifecycle status. & ewasteId, ownerId, categoryId, imageUrls, aiPredictedCategory, lifecycleStatus \\
pickup\_requests & Stores pickup schedules, locations, requester information, and request status. & pickupRequestId, requesterId, ewasteIds, preferredDate, latitude, status \\
assignments & Stores collector, driver, technician, and recycler assignments. & assignmentId, referenceId, assignedUserId, assignedRole, assignmentStatus \\
transportation\_records & Stores route, origin, destination, driver, vehicle, and delivery data. & transportId, driverId, originAddress, destinationAddress, routeDistance, status \\
repair\_records & Stores inspection, diagnosis, parts, costs, repair progress, and final outcome. & repairId, ewasteId, technicianId, diagnosis, repairStatus, finalOutcome \\
refurbishment\_records & Stores refurbishment activities, testing results, quality grade, and approval status. & refurbishmentId, ewasteId, activitiesPerformed, qualityGrade, approvalStatus \\
recycling\_records & Stores dismantling, hazardous material, residual waste, and recycling outcome information. & recyclingId, ewasteId, recyclerId, hazardousMaterials, disposalMethod, recyclingStatus \\
recovered\_materials & Stores materials and components recovered during recycling. & materialRecordId, recyclingId, materialName, quantity, unit, estimatedValue \\
marketplace\_listings & Stores refurbished devices, reusable components, and recovered materials listed for sale or enquiry. & listingId, sourceRecordId, title, condition, price, listingStatus \\
notifications & Stores targeted alerts, announcements, workflow notifications, and read status. & notificationId, recipientId, notificationType, title, isRead, createdAt \\
environmental\_reports & Stores environmental reports, observations, findings, and recommendations. & reportId, reportType, generatedBy, findings, recommendations, status \\
audit\_logs & Stores significant user, operational, security, and administrative activities. & auditId, actorId, action, referenceId, result, createdAt \\
system\_settings & Stores configurable non-secret system settings. & settingId, settingName, settingValue, category, updatedBy \\
reference\_data & Stores controlled lists such as categories, statuses, grades, facility types, and report types. & referenceId, referenceType, name, code, isActive \\
\bottomrule
\end{longtable}

\subsection{User Interface Design}

User interface design defines how users interact with the EcoTrace system
through mobile, web, and Windows applications. The interface design
translates the functional requirements into screens, navigation structures,
forms, dashboards, tables, maps, notifications, and interactive controls. The
EcoTrace user interfaces were designed to be clear, consistent, role-specific,
responsive, and accessible. Each actor is presented with functions relevant to
their responsibilities, while restricted operations remain hidden or
inaccessible. Flutter provides the shared presentation framework used to
implement the Android, Web, and Windows interfaces. Although the
applications share common components and design principles, their layouts
are adapted to different screen sizes and usage contexts.

\subsubsection{User Interface Design Objectives}
The primary objectives of the EcoTrace interface design are to: provide
simple and understandable navigation; present role-specific information and
operations; reduce the number of steps required to complete common tasks;
maintain visual consistency across mobile, web, and Windows applications;
provide clear validation, success, loading, and error feedback; support
different screen sizes and device orientations; make electronic waste
workflows easy to understand; present maps, charts, reports, and status
information clearly; protect restricted functionality through role-based
navigation; and support users with different levels of technical experience.

\subsubsection{Visual Design Principles}
The EcoTrace visual identity uses environmental and technology-related
design elements. Green is used to represent sustainability, recycling,
environmental protection, and successful operations. Blue is used to represent
technology, trust, security, and information management. The interface
design applies consistent colour usage, readable typography, sufficient
spacing between interface elements, clear visual hierarchy, recognizable
icons, consistent button styles, understandable status indicators, responsive
cards and panels, limited visual clutter, and a clear distinction between
primary and secondary actions. Red, amber, and green status colours may be
used carefully to represent errors, warnings, pending activities, and completed
operations; however, important information shall not depend on colour alone.

\subsubsection{Navigation Design}
EcoTrace uses role-based navigation to ensure that users access only the
functions relevant to their responsibilities. The mobile application may use a
bottom navigation bar, a navigation drawer, tab-based navigation, app-bar
actions, and floating action buttons for important operations. The Web and
Windows applications may use a collapsible side navigation panel, a top
navigation bar, breadcrumb navigation, page-level tabs, dashboard shortcuts,
and contextual action menus. The navigation system shall preserve the user's
current role and prevent unauthorized access through direct routes or
manually entered URLs.

\subsubsection{Responsive Layout Design}
Responsive design ensures that EcoTrace interfaces remain usable across
different screen sizes. The mobile interface uses a single-column layout for
most forms and content cards. The Web and Windows interfaces use wider
multi-column layouts, side panels, tables, charts, and map views where
suitable. Responsive behaviour includes flexible widths, adaptive grids,
scalable cards, scrollable tables, collapsible navigation, responsive charts and
maps, alternative mobile layouts for complex dashboards, and prevention of
horizontal overflow.

\subsubsection{Screen-by-Screen Interface Summary}

\textbf{Authentication interfaces} (registration, login, password recovery,
email verification, account status) -- the login interface contains an
email-address field, password field, show/hide password control, login
button, forgot-password link, account-registration link, loading indicator, and
authentication error message. Registration collects only the information
required to create an account and establish role/account type; detailed profile
information may be completed afterward.

\textbf{Household and Business dashboard} -- total registered items, active
pickup requests, completed collections, items under repair/recycling, recent
notifications, nearby collection centres, environmental awareness content,
quick actions, and recent lifecycle activity; business users additionally see
bulk-request, organizational-record, certificate, and report information.

\textbf{Electronic waste registration interface} -- category selector,
device-name field, brand/model fields, condition selector, quantity field,
description field, image-upload control, location selector, preferred recovery
option, and submit button, with validation messages and a confirmation +
generated record identifier on success.

\textbf{AI classification interface} -- image preview, replacement control,
classification button, processing indicator, predicted category, confidence
score (displayed without implying absolute certainty), accept/correct-result
controls, and error message.

\textbf{Pickup scheduling interface} -- registered-item selector,
pickup-address field, interactive map, latitude/longitude values, preferred
date/time, contact details, collection instructions, estimated pickup summary,
and submit button; status values include submitted, approved, assigned, en
route, collected, completed, cancelled, or delayed.

\textbf{Collector and Driver interfaces} -- the Collector dashboard shows
assigned pickups, accepted tasks, upcoming schedules, completed collections,
pickup locations, waste details, requester contact, evidence controls, and
status-update actions; the Driver dashboard shows assigned tasks, a route map,
origin/destination, estimated distance/time, navigation, delivery information,
status, and incident reporting.

\textbf{Technician interface} -- assigned devices, pending inspections,
repairs in progress, completed repairs, refurbishment tasks, device
images/history, required parts, cost information, test results, and
final-outcome controls.

\textbf{Recycler interface} -- assigned recycling items, waste received,
tasks in progress, recovered materials, hazardous materials, residual waste,
processing facility, status, and completed records, supporting multiple
recovered materials per recycling process.

\textbf{Environmental Officer dashboard} -- location map, coverage map,
active/delayed cases, collection/repair/recycling statistics, recovered-material
indicators, facility activity, alerts, inspection records, charts/trends, and
report generation, filterable by date, district, category, facility, status, or
actor.

\textbf{Administrator dashboard} -- registered users, pending verification,
registered waste, new requests, active assignments/collections, repair/recycling
workflows, marketplace listings, notifications/announcements, operational
charts, recent activity, and unresolved issues; tables support search, filtering,
sorting, pagination, and contextual actions.

\textbf{Super Administrator dashboard} -- administrator management,
roles/permissions, audit logs, security alerts, system settings, integration and
API configuration status, account suspensions, system-health indicators, and
high-risk action confirmations, with clear confirmation required for sensitive
operations.

\textbf{Marketplace interface} -- product cards/images, title/category,
condition/quality grade, price, availability, seller/facility information,
search/filters, product-detail screen, and enquiry/purchase action; only
approved and active listings are displayed publicly.

\textbf{Notification interface} -- title, message, type, date/time, read/unread
status, related-record link, and action button, with mark-as-read support.

\textbf{Reports and analytics interface} -- summary cards, tables, bar/line/pie
charts, map visualizations, date/category/status filters, and export controls,
showing only information permitted for the user's role.

\subsubsection{Form, Feedback, and Error-Handling Design}
Forms across EcoTrace shall follow a consistent structure: clear title, short
instructions where necessary, labelled fields, required-field indicators,
suitable input controls, validation messages, primary and secondary actions,
loading feedback, and success or error response; long forms should be divided
into logical sections or steps. Feedback mechanisms include confirmation
dialogs, snack-bar messages, notification banners, loading and progress
indicators, error messages, empty states, status badges, and completion
summaries, with status labels kept consistent across platforms. Error messages
should describe what went wrong, whether the user can correct the issue,
whether to retry, whether a service is temporarily unavailable, and whether
administrator support is required -- technical stack traces, credentials,
internal service information, and sensitive error details shall not be displayed
to standard users.

\subsubsection{Accessibility Considerations}
The interface design shall consider accessibility through readable font sizes,
suitable colour contrast, clear form labels, descriptive icons, sufficiently large
touch controls, logical navigation order, keyboard support for the Web
application, responsive layouts, and visible error and status information.

\subsubsection{User Interface Design Summary}
The EcoTrace user interface design provides role-specific access through
Android, Web, and Windows applications. The design supports electronic
waste registration, pickup scheduling, collection, transportation, repair,
recycling, monitoring, reporting, administration, and marketplace activities.
The use of consistent layouts, responsive components, clear navigation,
validation feedback, dashboards, maps, charts, and workflow statuses is
intended to improve usability and reduce operational errors.
Figure~\ref{fig:ui-framework} illustrates the role-based and cross-platform
user interface design structure adopted for EcoTrace, and
Table~\ref{tab:ui-summary} summarizes the major interfaces.

\begin{figure}[H]
  \centering
  \imageplaceholder{User Interface Design Framework for EcoTrace}
  \caption{User Interface Design Framework for EcoTrace. Source: Author's design.}
  \label{fig:ui-framework}
\end{figure}

\begin{longtable}{@{}p{3.6cm}p{3.4cm}p{6.2cm}@{}}
\caption{Summary of Major EcoTrace User Interfaces} \label{tab:ui-summary} \\
\toprule \textbf{Interface} & \textbf{Primary Users} & \textbf{Principal Functions} \\ \midrule
\endfirsthead
\toprule \textbf{Interface} & \textbf{Primary Users} & \textbf{Principal Functions} \\ \midrule
\endhead
Authentication interface & All users & Registration, login, password recovery, email verification, and account status feedback. \\
Household and Business dashboard & Household and Business Users & Waste registration, pickup requests, tracking, notifications, marketplace access, and personal or organizational reports. \\
Collector dashboard & Collectors & Pickup assignments, waste details, location access, collection status, evidence upload, and task history. \\
Driver dashboard & Drivers & Transportation assignments, routes, origin and destination information, status updates, incident reporting, and delivery confirmation. \\
Technician dashboard & Technicians & Inspection, diagnosis, repair, refurbishment, parts, testing, quality grading, and final outcomes. \\
Recycler dashboard & Recyclers & Waste receipt, dismantling, recovered materials, hazardous materials, residual waste, processing status, and recycling completion. \\
Environmental Officer dashboard & Environmental Officers & Maps, monitoring, collection coverage, environmental indicators, facility review, alerts, analytics, and reports. \\
Administrator dashboard & Administrators & User management, pickup review, assignments, facilities, categories, marketplace management, announcements, reports, and operational monitoring. \\
Super Administrator dashboard & Super Administrators & Administrator management, roles, permissions, audit logs, security, configuration, integrations, and system oversight. \\
Marketplace interface & Registered Users and Authorized Sellers & Browse listings, search, filter, view product details, manage approved listings, and submit enquiries. \\
Reporting interface & Business Users, Environmental Officers, Administrators, and Super Administrators & Tables, charts, maps, filters, summaries, analytics, and report export. \\
\bottomrule
\end{longtable}

\subsection{Hardware, Software, and Platform Requirements}

The successful development, deployment, and use of EcoTrace depend on
appropriate hardware, software, network, development, and cloud-service
requirements. These requirements define the minimum resources needed by
developers, administrators, operational personnel, and end users. The
requirements are grouped into client-device requirements, development
requirements, server and cloud requirements, network requirements,
external-service requirements, and platform-specific requirements.

\subsubsection{Client Hardware Requirements}
EcoTrace supports Android mobile devices, Web browsers, and Windows
desktop computers. The Android application should operate on a device with
a multi-core mobile processor, at least 2~GB RAM, at least 200~MB available
storage, a touchscreen display, a rear camera, GPS/location-service support,
and Wi-Fi or mobile-data connectivity; a device with at least 4~GB RAM, a
modern processor, reliable GPS, and a high-quality camera is recommended
for field roles. The Web application requires a dual-core or higher processor,
at least 4~GB RAM, a display of at least 1366$\times$768, keyboard and
pointing device, and a stable Internet connection. The Windows application
requires a 64-bit processor, at least 4~GB RAM, at least 500~MB storage, a
display of at least 1366$\times$768, keyboard/mouse or touch input, Internet
connectivity, and a supported Windows version; administrators and
environmental officers are recommended at least 8~GB RAM and a Full HD
display.

\textbf{Optional operational hardware} that may improve efficiency
includes QR-code or barcode scanners, digital weighing scales, label printers,
GPS-enabled collection vehicles, mobile receipt printers, external cameras,
rugged mobile devices, power banks, and protective equipment. Future
versions may integrate IoT sensors, smart bins, vehicle-tracking devices, and
weighing systems.

\subsubsection{Development Hardware Requirements}
The minimum development hardware includes a 64-bit multi-core
processor, at least 8~GB RAM, at least 20~GB free storage, reliable Internet
connectivity, a display of at least 1366$\times$768, and hardware
virtualization support where an Android emulator is used. The recommended
configuration is an Intel Core i5/AMD Ryzen 5 or equivalent processor,
16~GB or more RAM, solid-state storage with at least 50~GB free, a Full HD
display, hardware virtualization support, and graphics capable of supporting
the development tools. A GPU may improve PyTorch model training speed,
though training may also be performed on CPU or an external cloud
environment.

\subsubsection{Client Software Requirements}
The Android device requires a supported Android OS version, Google Play
Services where required by Firebase and Google Maps, camera/location/file/
notification permissions where applicable, current security updates, and the
installed EcoTrace application. The Web application requires a current
mainstream browser (Google Chrome, Microsoft Edge, Mozilla Firefox, or
another standards-compliant browser verified during testing) supporting
JavaScript, HTTPS, responsive layouts, and the Web APIs required by Flutter
Web. The Windows application requires a supported 64-bit Windows OS, the
installed EcoTrace application, required runtime components, and access to
the Internet and required cloud-service endpoints.

\subsubsection{Development Software Requirements}
The development environment includes the Flutter SDK, the Dart
programming language, Visual Studio Code or Android Studio, the Android
SDK, the Python programming language, FastAPI, the Uvicorn ASGI server,
PyTorch and TorchVision, Firebase command-line and configuration tools,
Git and GitHub, Web browser development tools, Postman or an equivalent
API-testing tool, LaTeX and Overleaf for documentation, and platform-specific
build tools for Windows and Android. \textit{The exact software versions used
in the final implementation are recorded in Chapter~\ref{ch:implementation},
since compatibility may change between Flutter, Dart, Python, PyTorch,
Firebase, and operating-system releases.}

\subsubsection{Backend and Server Software Requirements}
The EcoTrace backend requires a Python runtime, FastAPI, Uvicorn,
Pydantic, the Firebase Admin SDK, the Cloudinary Python SDK, Google Maps
service integration, Firebase Cloud Messaging integration, PyTorch
model-loading and inference dependencies, environment-variable
management, dependency specification through \texttt{requirements.txt}, and
secure HTTPS access through the deployment platform. The backend shall
expose RESTful API endpoints and exchange structured data using JSON.

\subsubsection{Cloud Platform Requirements}
\textbf{Render} hosts the FastAPI backend and provides Git-based
deployment, build/start commands, a public HTTPS endpoint, environment
variables, runtime and deployment logs, and health-check support; the free
tier is suitable for academic demonstration but may experience cold starts and
limited compute capacity. \textbf{Cloud Firestore} requires a configured
Firebase project, an approved database region, collections and indexes,
Firebase Security Rules, service-account access for the backend, and network
access to Firebase services. \textbf{Firebase Authentication} requires enabled
authentication providers, configured client applications, approved domains
for Web authentication, email/password authentication where used, and
secure token verification by the backend. \textbf{Firebase Cloud Messaging}
requires configured application identifiers, registered device/browser tokens,
notification permissions, backend messaging credentials, and supported
client-platform notification services. \textbf{Cloudinary} requires a configured
cloud account, cloud name, securely stored API key/secret, approved upload
settings, media-transformation configuration, and secure media URLs.
\textbf{Google Maps Platform} requires enabled Maps APIs/SDKs, restricted
API keys, an authorized Android package name and certificate where
applicable, authorized Web domains, billing/quota configuration, usage
monitoring, and network access to Google Maps services.

\subsubsection{Artificial Intelligence Requirements}
The AI component requires a labelled electronic waste image dataset, an
image-preprocessing pipeline, training/validation/test datasets, PyTorch and
TorchVision, a selected CNN or transfer-learning architecture, trained model
parameters, model evaluation metrics, model-loading logic within the
FastAPI backend, image validation and transformation, sufficient CPU/
memory/storage resources, and acceptable model size and inference time for
the deployment environment. The final model architecture was to be selected
after evaluating classification accuracy, precision, recall, F1-score, model
size, memory usage, and inference time.

\subsubsection{Network, Security, Storage, and Platform Compatibility
Requirements}
EcoTrace depends on stable Internet connectivity, HTTPS access to Firebase,
Cloudinary, Google Maps Platform, and the Render-hosted backend, and
sufficient upload/download speed for images, maps, reports, and dashboard
data; the system should provide loading states, retry options, and clear error
messages when network services are unavailable. Security-related platform
requirements include HTTPS for client-server communication, Firebase
Authentication token verification, role-based backend authorization, Firestore
Security Rules, restricted Google Maps API keys, environment variables for
backend secrets, secure Cloudinary credentials, controlled service-account
permissions, validation of uploaded files, secure error handling, audit
logging, and separation of development/production configuration where
possible -- secrets shall not be committed to the GitHub repository or
embedded in public client code. Structured data is stored in Cloud Firestore
and images in Cloudinary; the system shall avoid storing permanent records
only on local device/browser/desktop storage or the ephemeral Render
filesystem. The Android, Web, and Windows applications shall support
consistent core functionality (authentication, role-based navigation, waste
records, pickup requests, workflow status, notifications, reports, and secure
backend communication), while some platform-specific functions may differ
(e.g. direct camera/GPS access on Android, browser permissions on Web,
wider desktop-oriented layouts on Windows).

\subsubsection{Recommended Production Requirements}
For large-scale or public production deployment, recommended
improvements include upgrading from the Render Free tier, deploying
multiple backend instances where required, configuring monitoring and
alerting, establishing formal backup procedures, configuring API quotas and
budget alerts, using production-grade logging and error monitoring,
performing security and load/performance testing, establishing
disaster-recovery procedures, introducing automated CI/CD, reviewing
applicable privacy and environmental regulations, and providing operational
support and incident-response procedures.

\subsubsection{Requirements Summary}
The hardware, software, network, and platform requirements define the
technical environment necessary to develop, deploy, and operate EcoTrace.
The system can be accessed through Android mobile devices, Web browsers,
and Windows computers. Its backend and cloud functionality depend, as
proposed at this stage of the project, on FastAPI, Render, Cloud Firestore,
Firebase Authentication, Firebase Cloud Messaging, Cloudinary, Google Maps
Platform, and PyTorch. \textit{Chapter~\ref{ch:implementation} records the
exact devices, operating systems, package versions, development tools, cloud
configuration, and deployment environment actually used during system
development and testing.} Table~\ref{tab:hw-requirements} and
Table~\ref{tab:sw-requirements} summarize the minimum and recommended
requirements.

\begin{table}[H]
\centering
\small
\caption{Minimum and Recommended Client Hardware Requirements}
\label{tab:hw-requirements}
\begin{tabular}{@{}p{2.6cm}p{5.4cm}p{5.4cm}@{}}
\toprule
\textbf{Platform} & \textbf{Minimum Requirement} & \textbf{Recommended Requirement} \\
\midrule
Android Mobile & Multi-core processor, 2~GB RAM, camera, GPS, Internet connection, sufficient storage. & Modern processor, at least 4~GB RAM, reliable GPS, quality camera, stable mobile-data or Wi-Fi. \\
Web Client & Dual-core processor, 4~GB RAM, modern browser, 1366$\times$768 display, Internet connection. & Quad-core processor, at least 8~GB RAM, Full HD display, current browser, stable broadband. \\
Windows Desktop & 64-bit processor, 4~GB RAM, 500~MB free storage, 1366$\times$768 display, Internet connection. & Modern 64-bit processor, at least 8~GB RAM, SSD storage, Full HD display, stable broadband. \\
Development Computer & 64-bit multi-core processor, 8~GB RAM, 20~GB free storage, virtualization support. & Core i5/Ryzen 5 equivalent, 16~GB RAM, SSD storage, at least 50~GB free space, optional GPU. \\
\bottomrule
\end{tabular}
\end{table}

\begin{longtable}{@{}p{3cm}p{4cm}p{6.2cm}@{}}
\caption{Software and Platform Requirements for EcoTrace} \label{tab:sw-requirements} \\
\toprule \textbf{Category} & \textbf{Technology or Service} & \textbf{Purpose} \\ \midrule
\endfirsthead
\toprule \textbf{Category} & \textbf{Technology or Service} & \textbf{Purpose} \\ \midrule
\endhead
Client Development & Flutter and Dart & Development of Android, Web, and Windows applications from a shared codebase. \\
Backend Development & Python, FastAPI, Uvicorn, and Pydantic & Implementation of REST APIs, validation, business logic, and backend services. \\
Database & Cloud Firestore & Storage and synchronization of structured application data. \\
Authentication & Firebase Authentication & Registration, login, password recovery, identity tokens, and session management. \\
Notifications & Firebase Cloud Messaging & Delivery of targeted alerts, workflow updates, and announcements. \\
Media Hosting & Cloudinary & Storage, optimization, transformation, and delivery of images. \\
Mapping and Routing & Google Maps Platform & Maps, geocoding, locations, routes, distance calculation, and route optimization. \\
Backend Deployment & Render & Hosting, HTTPS, logging, environment variables, and Git-based deployment. \\
Artificial Intelligence & PyTorch and TorchVision & Training and deployment of the electronic waste image-classification model. \\
Source Control & Git and GitHub & Version control, collaboration, change tracking, and deployment integration. \\
Development Environment & Visual Studio Code or Android Studio & Editing, debugging, testing, and running Flutter and Python code. \\
API Testing & Postman or equivalent tool & Testing backend endpoints, requests, responses, authentication, and error handling. \\
Documentation & LaTeX and Overleaf & Preparation and formatting of the final project documentation. \\
\bottomrule
\end{longtable}

\section{Chapter Summary}

This chapter presented the analysis and design of the EcoTrace: Intelligent
Electronic Waste Recovery and Circular Economy Management System.

The chapter began by describing the adapted Agile software development
methodology used during the project. The methodology supported iterative
planning, analysis, design, implementation, testing, review, deployment, and
continuous refinement of the system modules.

The existing electronic waste management process was then analysed. The
analysis showed that the current process is largely manual, fragmented, and
dependent on informal communication and unstructured records. Major
problems identified included poor lifecycle traceability, inefficient pickup
scheduling, limited route planning, weak stakeholder coordination, unsafe
disposal practices, inadequate environmental monitoring, weak reporting,
limited security, and insufficient support for circular economy activities.

The proposed EcoTrace system was presented as an integrated,
cross-platform, and cloud-based solution intended to address these problems.
The system supports household users, business users, collectors, drivers,
technicians, recyclers, environmental officers, administrators, and super
administrators.

The functional requirements defined the services that EcoTrace shall
provide, including authentication, role-based access control, electronic waste
registration, AI-assisted classification, pickup scheduling, collection,
transportation, repair, refurbishment, recycling, marketplace management,
notifications, reporting, environmental monitoring, administration, audit
tracking, and external-service integration.

The non-functional requirements defined the expected quality attributes of
the system, including performance, availability, reliability, security, privacy,
usability, accessibility, scalability, maintainability, compatibility, data
integrity, backup, recovery, auditability, portability, external-service
resilience, and ethical operation.

The system design described the layered architecture of EcoTrace, as
proposed at this stage of the project: Flutter provides the Android, Web, and
Windows presentation layers, while FastAPI provides the central application
and business-logic layer, deployed on Render and communicating with Cloud
Firestore, Firebase Authentication, Firebase Cloud Messaging, Cloudinary,
Google Maps Platform, and a PyTorch artificial intelligence component. The
use case model identified the responsibilities and interactions of the major
actors; the context-level and Level-One Data Flow Diagrams described the
movement of information between users, processes, data stores, and external
services; and sequence diagrams illustrated the order of interactions involved
in authentication, waste registration, AI classification, pickup assignment,
logistics, recovery, and reporting workflows.

The database design presented the major Cloud Firestore collections and
their logical relationships, supporting user profiles, electronic waste records,
pickup requests, assignments, transportation, repair, refurbishment, recycling,
recovered materials, marketplace listings, notifications, environmental
reports, audit logs, settings, and reference data.

The user interface design established a shared Flutter design system with
responsive, role-based interfaces for Android, Web, and Windows. The
chapter also specified the hardware, software, network, cloud, security, and
platform requirements necessary for development, deployment, and
operation.

The analysis and design presented in this chapter provide the technical
foundation for the implementation, testing, results, and evaluation described
in Chapter~\ref{ch:implementation} -- which, per NFR-ETH-9
(Section~\ref{sec:nfr}), reports plainly where the delivered system's
technology stack and functional coverage followed this chapter's proposal and
where it took a different, equally deliberate path.
